\documentclass[]{Liebert_Author_revised}
%%Liebert_Author_revised.cls and plainnat-revised.bst edited by Tandy Warnow
\usepackage{graphicx}
\usepackage{multicol}
\usepackage{multirow}

\usepackage{url}
\usepackage{color}
\usepackage{marvosym}
\usepackage{float}

\usepackage{caption}
\usepackage{subcaption}
\usepackage{cite}
\usepackage{amsmath,amssymb,amsfonts}
\usepackage{algorithmic}
\newcommand{\ie}{\emph{i.e.}}
\newcommand{\eg}{\emph{e.g.}}


% \title{Uncovering Hierarchical Structure in LLM Embeddings Using Geometric and Topological Measures} 

\title{
Uncovering Hierarchical Structure in LLM Embeddings Using Geometric and Topological Measures
}

% \author{Author 1,$^{1\ast}$ Author 2,$^{1}$ Author 3$^{2}$\\
% {$^{1}$Department of Chemistry, University of Wherever,}\\
% {An Unknown Address, Wherever, ST 00000, USA}\\
% {$^{2}$Another Unknown Address, Palookaville, ST 99999, USA}\\
% {$^\ast$To whom correspondence should be addressed;}\\
% {E-mail:  jsmith@wherever.edu.}
% }

\author{Prakash Chourasia$^1$, Sarwan Ali$^2$, Murray Patterson$^1$, and Alex Zelikovsky$^1$\\
  {$^1$Georgia State University} 
  {Atlanta, GA 30303, USA}\\
  {$^2$Columbia University, New York, NY, USA}\\
E-mail: pchourasia1@student.gsu.edu, sa4559@cumc.columbia.edu, \\
   mpatterson30@gsu.edu, azelikovsky@cs.gsu.edu \\ 
}


\date{}

\begin{document} 

\maketitle 

\keywords{Large Language Models (LLMs), Geometric Group Theory, Embedding Analysis, Representation Learning, Metric Space Analysis}

\begin{abstract}
The rapid advancement of large language models (LLMs) has driven major advances across scientific domains, yet the hierarchical, geometric, and topological structure of their embedding spaces remains poorly understood. In this work, we introduce a comprehensive framework for evaluating LLM embeddings through six complementary metrics that characterize hierarchy, geometry, and topology. Specifically, $\delta$-hyperbolicity and ultrametricity quantify the degree of tree-like structure, while Neighbor Joining evaluates how well pairwise distances support hierarchical reconstruction. The Ollivier-Ricci curvature captures local geometric behavior, persistent homology measures multiscale topological stability, and fractal dimension estimates intrinsic complexity. Applying this framework to embeddings from multiple protein LLMs, we reveal how different models encode hierarchical relationships and structural organization. Evaluating the structure of these embeddings in the representation space helps improve geometry-aware model design. This is critical in domains such as bioinformatics and healthcare, where real-world hierarchical and relational structures matter. Our results demonstrate that geometric and topological analysis provides a guiding lens for understanding and potentially improving the representational quality of LLM embeddings.
\end{abstract}



\section{Introduction}
\label{sec_intro}
The integration of large language models (LLMs) in healthcare has advanced rapidly. These models generate embeddings that capture rich contextual and sequential information, enabling downstream tasks such as protein function prediction and structure inference. This mapping of proteins in hyperbolic space~\citeauthor{dress2012basic} enables efficient clustering of functionally similar proteins, thus supporting applications in Functional Annotation of Proteins~\citeauthor{ali2023protein}. However, understanding the geometric properties of the embeddings is crucial to improving reliability~\citeauthor{nickel2017poincare}.
One way to assess the structural properties of embeddings is through $\delta$-hyperbolicity, a concept introduced by Gromov~\citeauthor{gromov1987hyperbolic}. $\delta$-hyperbolicity quantifies the extent to which a metric space deviates from being an ideal tree~\citeauthor{yang2024enhancing} like structure, where lower $\delta$ values indicate structures closer to a tree-like geometry, which suggests hierarchical relationships among data points. 
In contrast, a higher  $\delta$ indicates a more complex, non-hierarchical structure. 
A major empirical study tested whether social and information networks are tree-like using Gromov’s $\delta$-hyperbolicity~\citeauthor{adcock2013tree}, showing that although traditional metrics miss such structure, hyperbolic analysis reveals meaningful hierarchical patterns.
$\delta$-hyperbolicity quantifies approximate tree-likeness in complex biological networks, simplifying tasks such as distance estimation, network classification, and core–periphery identification~\citeauthor{abu2016metric}. Ultrametricity~\citeauthor{chierchia2019ultrametric} is another popular stricter metric for capturing hierarchical structure, ideal for deterministic, nested data such as phylogenetic trees. In techniques like UPGMA (Unweighted Pair Group Method with Arithmetic Mean) and WPGMA (Weighted Pair Group with Arithmetic Mean), which are used to build ultrametric trees in hierarchical clustering, the idea of ultrametricity is fundamental~\citeauthor{dubes1999cluster}. Whereas Q-matrix statistics from Neighbor Joining offer a quantitative, algorithmic evaluation of tree-likeness that complements these geometric metrics. Strong adherence to additive tree structure is indicated by low normalized Q values~\citeauthor{gascuel2006neighbor,atteson1999performance}. NJ evaluates whether pairwise distances are compatible with an additive tree by transforming the distance matrix into a Q-matrix, demonstrating how strongly each pair of points supports hierarchical merging. 

Beyond tree-likeness, embedding geometry can exhibit additional structural properties that strictly hierarchical metrics do not capture. We incorporate three complementary geometric and topological measures that probe different, orthogonal aspects of the space. These include Ollivier–Ricci curvature~\citeauthor{ollivier2009ricci, ni2015ricci} for local geometric behavior, persistent homology~\citeauthor{edelsbrunner2008persistent} for multiscale topological structure, and fractal dimension~\citeauthor{grassberger1983measuring} for intrinsic geometric complexity. 

We adapt $\delta$-hyperbolicity, ultrametricity, neighborhood joining, and three complementary geometric and topological measures to the domain of genomic sequencing data analyses by evaluating the embeddings generated by LLMs such as ``Tasks Assessing Protein Embeddings (TAPE)"~\citeauthor{rao2019evaluating}, ESM2~\citeauthor{lin2023evolutionary}, Seqvec~\citeauthor{heinzinger2019modeling}, and ProtT5~\citeauthor{elnaggar2021prottrans}. We aim to provide a deeper understanding of the geometric properties underlying LLM-generated representations.
% The key contributions of this study are as follows:
% \begin{enumerate}
%     \item We generate embeddings for three datasets using various large language models (LLMs), and evaluate them using $\delta$-hyperbolicity and ultrametricity to assess the geometric properties of the embedding spaces, and comparative analysis of embeddings generated by several LLMs.
%     \item We perform classification on the embeddings to further evaluate their quality and provide empirical validation for the two geometric metrics.
%     \item We also perform clustering on the embeddings to evaluate their quality and provide empirical validation.
%     % \item We visualize the embeddings using t-SNE plots to gain qualitative insights into the structure and separability of the representations by embeddings generated by several LLMs.
% \end{enumerate}
The key contributions of this study are as follows:
\begin{enumerate}
\item We generate embeddings for three datasets using different LLMs and evaluate them with six complementary geometric and topological metrics spanning hierarchy, local geometry, topology, and intrinsic complexity.
\item We provide a comparative geometric analysis of the embeddings of several protein LLMs (SeqVec, ESM2, TAPE, ProtT5), showing how the model design shapes the embedding structure.
\item We assess embedding quality through classification and clustering, demonstrating that geometric properties correlate with downstream performance.
\end{enumerate}

   
% In the remaining sections, we provide a brief literature review in Section~\ref{sec_rw}. The proposed approach and methodology are discussed in Section~\ref{sec_pa}. Section~\ref{sec_evaluation_metrics} discusses the metrics used for assessing the geometric properties. The details regarding the dataset statistics and experimental setup are presented in Section~\ref{sec_es}. We then present experimental results and comparative analyses with existing methods in Section~\ref{sec_results}. Finally, the conclusion of the whole study is presented in Section~\ref{sec_conclusion}.

\section{Related Work}
\label{sec_rw}

% We divide the literature work into several subcategories and discuss each one separately below.

% \subsection{Health Recommender Systems}
% Health recommender systems are designed to provide personalized health advice and interventions based on user data. Previous research has focused on developing algorithms that improve the precision and relevance of recommendations by incorporating various data sources, including electronic health records and patient self-reports~\citeauthor{lattar2020health, ccelik2020survey}. These systems have been applied to a wide range of health domains, from nutrition and fitness to chronic disease management~\citeauthor{hors2018analyzing}.

% \subsection{Large Language Models and Embeddings}
LLMs have become widely used to produce rich embeddings that boost performance across diverse tasks. Similar models are now used in protein biology to map sequence and structure for better function prediction~\citeauthor{ali2023protein}.
% \subsection{Hyperbolicity and Metric Space Analysis}
$\delta$-hyperbolicity, introduced by Gromov~\citeauthor{gromov1987hyperbolic}, provides a geometric measure of how tree-like a metric space is. It has been widely applied in graph theory and computational geometry~\citeauthor{krioukov2010hyperbolic} to analyze structural properties relevant to computational problems.
Authors in~\citeauthor{chami2020trees} show that H\textsc{yp}HC closely approximates the optimal tree and outperforms other clustering methods.
% Authors in~\citeauthor{chami2020trees} show that their method, Hyperbolic Hierarchical Clustering H\textsc{yp}HC, can closely match the best possible tree within a small error and often outperforms traditional and other learning-based clustering methods. 
% Instead of being defined in terms of the cost function that needs to be optimized, standard hierarchical clustering techniques are described procedurally. 
Authors in~\citeauthor{chierchia2019ultrametric} frame hierarchical clustering as an optimization problem, seeking the ultrametric that best fits the data. 
% They simplify the constraints by encoding them as a min–max term, enabling efficient gradient-based optimization.
% \subsection{Applications to Embeddings}

Recent work shows that geometric properties of embeddings, including hyperbolicity, can inform clustering and classification~\citeauthor{yang2024enhancing}, but these ideas have not been applied to biological sequence embeddings. Since many biological processes follow tree-like evolutionary patterns~\citeauthor{felsenstein2004inferring, abu2016metric}, $\delta$-hyperbolicity provides a natural tool for assessing and improving models in phylogenetics, viral evolution, and epidemiology. Hyperbolic embedding techniques have also shown benefits in drug–target prediction and functional protein annotation~\citeauthor{dress2012basic, ali2023protein}, highlighting the broader potential of hyperbolic geometry in biological modeling.
% Recent work has examined the geometry of high-dimensional embeddings, including hyperbolicity measures~\citeauthor{yang2024enhancing}, showing that geometric structure can inform tasks like clustering and classification. However, hyperbolicity has not yet been applied to biological sequence embeddings. Because many biological processes follow tree-like evolutionary relationships~\citeauthor{felsenstein2004inferring, abu2016metric}, $\delta$-hyperbolicity offers a promising way to assess and improve models used in phylogenetics, viral evolution, and epidemiology.
% Using hyperbolic embeddings guided by $\delta$-hyperbolicity improves drug-target interaction predictions. The mapping of proteins in a hyperbolic space~\citeauthor{dress2012basic} allows efficient clustering of functionally similar proteins, thus providing applications in Functional Annotation of Proteins~\citeauthor{ali2023protein}.
% \subsection{Geometric and Topological Analysis of Embeddings}
% \textcolor{green}{
Recent studies use geometric tools such as Ricci curvature, persistent homology, and fractal dimension to analyze neural representations~\citeauthor{ni2015ricci, carlsson2009topology, facco2017estimating}. We incorporate these ideas to provide a unified geometric and topological assessment of protein LLM embeddings.
% Some other recent work has explored geometric properties of neural network representations beyond traditional metrics. Ricci curvature has been applied to analyze graph neural networks~\citeauthor{ni2015ricci} and detect anomalies in network embeddings~\citeauthor{sandhu2015graph}.  Persistent homology provides a principled framework for multi-scale topological analysis~\citeauthor{carlsson2009topology, edelsbrunner2008persistent}, with applications to deep learning, 
% analyzing representational geometry~\citeauthor{naitzat2020topology}. Fractal-dimension–based intrinsic dimension estimation has helped reveal the geometry and low-dimensional structure of deep embeddings~\citeauthor{facco2017estimating}. Here, we integrate these geometric and topological tools to provide a comprehensive characterization of protein LLM embeddings.
% }


\section{Proposed Approach}
\label{sec_PAa}

In this section, we outline our framework for evaluating embeddings generated by large language models (LLMs). Our analysis proceeds as follows :
\begin{enumerate}
    \item Generate embeddings using the different LLMs.
    \item Compute Euclidean and Poincaré distance matrices.
    \item Evaluate tree-like structure using three complementary metrics: $\delta$-hyperbolicity, ultrametricity, and Neighbor Joining (via the distribution of Q-matrix).
    \item Assess additional geometric and topological structure using three complementary metrics.
    \item Analyze the classification and clustering results to understand the geometric properties of the embeddings.
\end{enumerate}

\subsection{LLMs for embedding generation}
We employ four LLMs, including SeqVec, ESM-2, TAPE, and ProtT5, to generate embeddings and evaluate how well they capture structural patterns in biological sequences.
\paragraph{Seqvec~\citeauthor{heinzinger2019modeling}}
SeqVec is an ELMo-inspired protein language model that learns contextual embeddings directly from large unlabeled protein datasets such as UniRef50~\citeauthor{suzek2015uniref}. It generates fast, alignment-free representations that capture meaningful sequence features and can rival traditional evolutionary methods~\citeauthor{sarzynska2021detecting}.
% SeqVec is a deep learning-based method that represents protein sequences as continuous vector embeddings without relying on evolutionary information. Inspired by the ELMo language model from NLP~\citeauthor{sarzynska2021detecting}, SeqVec treats amino acids as words and learns contextual relationships from large unlabeled protein databases like UniRef50~\citeauthor{suzek2015uniref}. Once trained, it generates rich, informative embeddings for new sequences in a fraction of a second, bypassing the time-consuming need for multiple sequence alignments. SeqVec demonstrates that single-sequence models can match or even outperform some alignment-based methods, offering a fast and scalable approach for protein function and structure prediction.

\paragraph{ESM2~\citeauthor{lin2023evolutionary}}
ESM-2 is a large transformer-based protein language model~\citeauthor{lin2023evolutionary} that predicts structural features directly from single sequences, eliminating the need for MSAs. Scaled to billions of parameters, it captures detailed structural signals~\citeauthor{rao2020transformer} and powers fast, high-accuracy tools like ESMFold, advancing single-sequence structure prediction.
% ESM-2 is a state-of-the-art protein language model developed by~\citeauthor{lin2023evolutionary}. It enables high-resolution protein structure prediction directly from single amino acid sequences, without the need for multiple sequence alignments (MSAs). Built on transformer architecture and scaled up to 15 billion parameters, ESM-2 captures detailed atomic-level structural information as an emergent property of the learned sequence representations~\citeauthor{rao2020transformer}. This model powers ESMFold, a fast and accurate structure prediction tool that approaches the performance of traditional MSA-based methods like AlphaFold, while being significantly faster—enabling predictions. ESM-2 represents a major leap forward in single-sequence-based structure prediction and has broad applications in structural biology, metagenomics, and protein function discovery.

\paragraph{TAPE~\citeauthor{rao2019evaluating}}
TAPE is a protein representation learning framework that trains transformer, LSTM, and ResNet models on large unlabeled datasets like UniRef50 using self-supervised objectives. It produces contextual embeddings for each sequence and provides the TAPE benchmark, which standardizes downstream tasks to consistently evaluate how well these representations capture structural and functional properties of proteins.
% TAPE is a semi-supervised LLM, protein representation learning method that works by training a protein large language model and then generating numerical embeddings. 
% The TAPE framework is built around a transfer learning approach for protein sequences
% % ~\citeauthor{rives2021biological}
% , using deep neural network architectures such as Transformers
% % ~\citeauthor{vaswani2017attention}
% , LSTMs
% % ~\citeauthor{hochreiter1997long}
% , and ResNets
% % ~\citeauthor{yu2017dilated}
% . These models are first trained on large, unlabeled protein sequence datasets, such as UniRef50~\citeauthor{suzek2015uniref}, using self-supervised learning objectives, notably masked language modeling, similar to BERT in NLP. The goal is to learn rich, general-purpose embeddings of protein sequences that can be fine-tuned or evaluated on downstream biological tasks, such as secondary-structure prediction, contact-map prediction, and protein stability. The TAPE benchmark standardizes these tasks, enabling consistent evaluation of how well learned representations capture structural and functional properties of proteins. We get the embeddings as an output, which reflects the contextualized representation of sequences after processing through the model.

\paragraph{ProtT5~\citeauthor{elnaggar2007prottrans}}
ProtT5 adapts the T5 transformer to protein sequences, training with a BERT-style masked language modeling objective and using only the encoder for efficient inference~\citeauthor{elnaggar2007prottrans}. Without requiring evolutionary data or MSAs, its embeddings capture significant biophysical characteristics and perform well on tasks such as secondary structure and localization prediction.
% ProtT5 is a powerful protein language model (pLM) based on the T5 (Text-to-Text Transfer Transformer) architecture from NLP, adapted to learn the language of proteins from large-scale amino acid sequence data. Unlike the original T5 model, which has both an encoder and decoder~\citeauthor{vaswani2017attention}. ProtT5 uses a BERT-style masked language modeling objective, focusing on reconstructing individual masked amino acids rather than spans. In ProtT5, only the encoder is used during inference because it performs better and is more efficient~\citeauthor{elnaggar2007prottrans}. ProtT5 embeddings capture key biophysical features of proteins, enabling accurate prediction of secondary structure, subcellular localization, and membrane association—all without relying on evolutionary information or multiple sequence alignments. Notably, ProtT5 achieved state-of-the-art accuracy in these tasks, suggesting it effectively learns protein grammar from raw sequence data alone.
% ProtT5 works by training a transformer encoder to reconstruct masked amino acids in protein sequences, thereby learning deep representations of protein "language". These embeddings enable accurate and efficient predictions of protein structure and function—without alignments, evolutionary data, or explicit labels during pretraining.





\subsection{Distance Matrix Calculation}

We compute pairwise distance matrices for all embeddings using both Euclidean and Poincaré metrics to obtain $N \times N$ distance matrices, where $N$ is the total number of protein sequences. Let $X$ be the set of embeddings.  

For the Euclidean metric, the distance between two embeddings $x_i$ and $x_j$ is given by:
\begin{equation}
\scriptsize
d_E(x_i, x_j) = \sqrt{\sum_{k=1}^d (x_{i,k} - x_{j,k})^2},
\end{equation}
where $d$ is the embedding dimensionality.  

For the Poincaré metric~\citeauthor{nickel2017poincare}, the hyperbolic distance between two sequence embeddings $x_i$ and $x_j$ within the unit ball is computed as:
\begin{equation}
\scriptsize
d_P(x_i, x_j) = \mathrm{arcosh} \left( 1 + 2 \frac{\|x_i - x_j\|^2}{(1 - \|x_i\|^2)(1 - \|x_j\|^2)} \right).
\end{equation}
Both metrics capture complementary geometric properties. Euclidean measures standard linear distances. Poincaré captures hierarchical, tree-like relationships in hyperbolic space.


% We compute the pairwise Euclidean distance matrix \(D\) for all pairs of embeddings to get an $N \times N$ dimensional distance matrix, where $N$ is the total number of protein sequences in our data. Let \(X\) be the set of embeddings where each row corresponds to an embedding generated from TAPE LLM. The Euclidean distance between two points \(x_i\) and \(x_j\) in \(X\) is computed as:

% \begin{equation}
% \scriptsize
% d(x_i, x_j) = \sqrt{\sum_{k=1}^d (x_{i,k} - x_{j,k})^2}
% \end{equation}

% where \(d\) is the dimensionality of the embeddings.



\section{Geometric and Topological Measures}
\label{sec_evaluation_metrics}

We evaluate embedding structure using six metrics spanning hierarchy, geometry, topology, and complexity: $\delta$-hyperbolicity, ultrametricity, Neighbor-Joining, Ollivier–Ricci curvature, persistent homology, and fractal dimension.

% We evaluate the tree-likeness using three metrics: $\delta$-hyperbolicity, ultrametricity, and Neighbor Joining (NJ). 

% $delta$-hyperbolicity (Gromov hyperbolicity) is a more flexible, approximate notion of tree-likeness. It applies to metric spaces and especially to graphs, where triangles are allowed to be “$delta$-slim,” meaning each side lies within a $delta$-radius of the other two sides. This permits local deviations or shortcuts while still maintaining an overall tree-like shape. It can be applied to more complex, noisy, or interconnected systems.

% Ultrametricity is a strong geometric condition where all triangles in the space are either isosceles with equal long sides or equilateral. It enforces a strict hierarchy, making it ideal for modeling perfect tree structures like phylogenetic trees and hierarchical clustering dendrograms. In essence, ultrametricity describes exact hierarchical tree. It can be used in deterministic, nested data (like biological taxonomies).

% Finally, we use Q-matrix statistics derived from the Neighbor Joining algorithm to quantify tree-likeness. Pairs of points with low (more negative) Q values indicate strong adherence to the NJ joining criterion, reflecting how well the distance matrix approximates an additive tree structure. Neighbor Joining (NJ) provides a quantitative, algorithmic perspective: by analyzing Q-matrix values, we can compute an NJ score that reflects how well the distance matrix approximates an additive tree. 


% Our use of Q-matrix statistics as a quantitative indicator of tree-likeness is motivated by theoretical analyses of Neighbor Joining that characterize NJ behavior directly in terms of distance-matrix properties~\citeauthor{gascuel2006neighbor,atteson1999performance}.

% Ultrametricity and $delta$-hyperbolicity both describe how closely a metric space or structure resembles a tree, but they differ significantly in rigidity and application. Ultrametricity is a strong geometric condition where all triangles in the space are either isosceles with equal long sides or equilateral. It enforces a strict hierarchy, making it ideal for modeling perfect tree structures like phylogenetic trees and hierarchical clustering dendrograms. 
% % In ultrametric spaces, the distance between any two points is no greater than the maximum of the other two pairwise distances, reflecting a tightly nested structure.
% In contrast, $delta$-hyperbolicity (Gromov hyperbolicity) is a more flexible, approximate notion of tree-likeness. It applies to metric spaces and especially to graphs, where triangles are allowed to be “$delta$-slim,” meaning each side lies within a $delta$-radius of the other two sides. This permits local deviations or shortcuts while still maintaining an overall tree-like shape. 
% % $delta$-hyperbolicity is particularly useful in analyzing large-scale networks, social graphs, and internet topologies where strict hierarchy doesn’t hold, but the structure still exhibits a tree-like tendency.

% In essence, ultrametricity describes exact hierarchical trees, while $delta$-hyperbolicity describes loosely tree-structured graphs. The former is used for deterministic, nested data (such as biological taxonomies), while the latter is applied to more complex, noisy, or interconnected systems.


\subsection{$\delta$-hyperbolicity}
The proposed algorithm for evaluating $\delta$-hyperbolicity is outlined in Algorithm~\ref{alg_hyperbolicity}  and is implemented to assess the geometric structure of LLM-derived protein sequence embeddings.
The algorithm begins by taking as input a set of embeddings \(X\) computed from different LLMs, a distance function \(d\), and a user-defined number of samples to evaluate. The goal is to calculate the $\delta$-hyperbolicity value \(\delta\) for the metric space defined by the embeddings. Initially, the embeddings are flattened and padded to a uniform length, then optionally reduced in dimensionality via PCA to retain 99\% of variance, ensuring efficient computation.
First, it computes the pairwise distance matrix \(D\) using the distance function \(d\) over the embeddings, where each entry in \(D\) represents the distance between a pair of embeddings. 
The core of the algorithm involves sampling a specified number of random quadruples \((a, b, c, w)\) from the embedding space, rather than iterating exhaustively over all combinations, which helps it to scale to large datasets.  
% For each quadruple, it computes the Gromov products \([a, b]_w\), \([b, c]_w\), and \([a, c]_w\) using the formula \([a, b]_w = \frac{1}{2} \left(d(a, w) + d(b, w) - d(a, b)\right)\), where \(d\) is the Euclidean or Poincare distance between embeddings.
% Next, it calculates the temporary slack term \(\delta_{\text{temp}}\) by taking the minimum of the Gromov products \([a, b]_w\) and \([b, c]_w\) and subtracting from \([a, c]_w\). This value represents the extent to which the space deviates from being hyperbolic (tree-like). After all quadruples have been evaluated, the algorithm returns \(\delta_{\text{max}}\), \(\delta_{\text{avg}}\), and \(\delta_{\text{std}}\) as the final $\delta$-hyperbolicity values, which quantifies the hyperbolicity of the metric space defined by the embeddings. This value provides insights into the geometric structure of the embeddings, which can be critical for understanding their effectiveness in underlying ML tasks.
To evaluate $\delta$-hyperbolicity, we need to compute the Gromov products for four points \(a\), \(b\), \(c\), and \(w\). The Gromov product \([a, b]_w\) is defined as (also mentioned in~\citeauthor{yang2024enhancing}):

\begin{equation}
\label{gromov_dist_eq1}
\scriptsize
[a, b]_w = \frac{1}{2} \left(d(a, w) + d(b, w) - d(a, b)\right)
\end{equation}

Where \(d\) is the Euclidean or Poincaré distance between embeddings.

The $\delta$-hyperbolicity is then calculated by checking if the following inequality holds for all sampled quadruples \(a\), \(b\), \(c\), and \(w\)~\citeauthor{yang2024enhancing}:

\begin{equation}
\label{gromov_dist_eq2}
\scriptsize
[a, c]_w \geq \min([a, b]_w, [b, c]_w) - \delta
\end{equation}

We estimate $\delta$-hyperbolicity by computing the minimum value of \(\delta\) that satisfies the above condition for all quadruples in the space.

To estimate the $\delta$-hyperbolicity of LLM-derived embeddings, we sample random quadruples of points in the embedding space and compute how much each quadruple deviates from the ideal tree-like geometry. After flattening and optionally reducing the embeddings with PCA, we compute pairwise distances and evaluate the Gromov products for each quadruple. The slack between these products gives a local estimate of the hyperbolic deviation. Aggregating these values across all samples yields the maximum, mean, and standard deviation of $\delta$, which together quantify how tree-like the embedding space is, where lower values indicate stronger hierarchical structure. The algorithm for computing $\delta$ Hyperbolicity is discussed in detail in the Supplementary Section.

Using parallel processing, these $\delta$ values are calculated for all quadruples sampled, after which the algorithm returns the maximum, average, and standard deviation of $\delta$ in all samples. These values collectively quantify the degree of hyperbolicity in the embedding space. A lower average $\delta$ indicates a more tree-like space, suggesting that the embedding structure may support hierarchical relationships.


\begin{algorithm}[h]
\caption{Compute $\delta$-Hyperbolicity}
\label{alg_hyperbolicity}
% \scriptsize
\small
\noindent
\begin{tabular}{p{1cm}p{6.8cm}}
\textbf{Input:}  & Set of embeddings \( X \), distance function \( d \), number of sampled quadruples \( S \) \\
\textbf{Output:} & Estimated $\delta$-hyperbolicity values: \( \delta_{\max}, \delta_{\text{avg}}, \delta_{\text{std}} \)
\end{tabular}

\begin{algorithmic}[1]

\STATE Sample \( S \) random quadruples \( (a, b, c, w) \) from \( X \)
\STATE Initialize empty list \( \Delta \leftarrow [\,] \)

% \STATE Compute pairwise distance matrix \(D\) from \(X\)
% \STATE Initialize \(\delta_{\text{min}} \leftarrow \infty\)

\FOR{each quadruple \( (a, b, c, w) \)}
    \STATE Compute Gromov products: \COMMENT{using Eq~\ref{gromov_dist_eq1} and Eq~\ref{gromov_dist_eq2}}
    \[
    [a, b]_w = \frac{1}{2} \left( d(a, w) + d(b, w) - d(a, b) \right)
    \]
    \[
    [b, c]_w = \frac{1}{2} \left( d(b, w) + d(c, w) - d(b, c) \right)
    \]
    \[
    [a, c]_w = \frac{1}{2} \left( d(a, w) + d(c, w) - d(a, c) \right)
    \]

     \STATE Compute temporary slack:
    \[
    \delta_{\text{temp}} = [a, c]_w - \min\left( [a, b]_w, [b, c]_w \right)
    \]

    % \STATE Compute slack term \(\delta\) using:
    % \[
    % \delta_{\text{temp}} = \max \left([a, b]_w, [b, c]_w\right) - [a, c]_w
    % \]
    % \STATE Update \(\delta_{\text{min}}\) with:
    % \[
    % \delta_{\text{min}} \leftarrow \min(\delta_{\text{min}}, \delta_{\text{temp}})
    % \]

    \STATE Append \( \delta_{\text{temp}} \) to \( \Delta \)
\ENDFOR

\STATE Compute:
\[
\delta_{\max} \leftarrow \max(\Delta), \]
% \quad 
% \\
\[
\delta_{\text{avg}} \leftarrow \text{mean}(\Delta), \]
% \quad 
% \\
\[
\delta_{\text{std}} \leftarrow \text{std}(\Delta)
\]

\RETURN \( \delta_{\max}, \delta_{\text{avg}}, \delta_{\text{std}} \)

\end{algorithmic}
\end{algorithm}


% \subsection{Hyperbolicity}


\subsection{Ultrametricity}

A metric space is said to be \textbf{ultrametric} if the distance function $d$ satisfies the \textit{ultrametric inequality}, also called a stricter version of the triangle inequality~\citeauthor{chierchia2019ultrametric}:

\begin{equation}
    \scriptsize
    d(i, k) \leq \max\{d(i, j), d(j, k)\} \quad \forall i, j, k
\end{equation}

This is a stronger condition than the standard triangle inequality. It enforces a strict hierarchy, making it ideal for modeling perfect tree structures like phylogenetic trees and hierarchical clustering dendrograms.


The algorithm~\ref{alg_ultrametricity_violations} determines whether a given distance matrix \( D \in \mathbb{R}^{n \times n} \) satisfies the ultrametric inequality. Specifically, the ultrametric condition requires that for any triplet of points \( (i, j, k) \), the largest pairwise distance among them must not exceed the maximum of the other two distances. To verify this, the algorithm iterates over all unique sampled combinations of triplets in the distance matrix. For each triplet, it retrieves the corresponding distances \( d_{ij} \), \( d_{ik} \), and \( d_{jk} \), then identifies the largest of the three. It compares this maximum value to the sum of the other two (smaller) distances, adjusted by a small tolerance \( \epsilon \) to account for numerical imprecision. If the maximum distance exceeds this sum by more than \( \epsilon \), the matrix is deemed to violate the ultrametric property, and the algorithm returns \textbf{False} immediately. If all triplets conform to the inequality, it concludes by returning \textbf{True}, confirming that the entire matrix exhibits ultrametricity. This algorithm is particularly valuable in fields such as phylogenetics and hierarchical clustering, where validating the ultrametric nature of data helps ensure the reliability of tree-based representations.

For each triplet, we check whether the largest of the three distances exceeds the other two by more than a small numerical tolerance. A violation is decided on the basis of it. We quantify the degree to which the distance matrix resembles a completely hierarchical ultrametric space by adding the maximum, mean, and standard deviation of these violations across all samples. As anticipated in applications such as phylogenetics and hierarchical clustering, this offers a straightforward method to assess whether the embeddings represent nested or tree-like relationships.

\begin{algorithm}[h]
\caption{Ultrametricity of a Distance Matrix}
\label{alg_ultrametricity_violations}
\noindent
% \scriptsize
\small
\begin{tabular}{p{1cm}p{7cm}}
\textbf{Input:}  & Distance matrix \(D \in \mathbb{R}^{n \times n}\), sample size \(S\), tolerance \(\epsilon\) \\
\textbf{Output:} & Dictionary containing ultrametricity violation metrics
\end{tabular}

\begin{algorithmic}[1]{}

\STATE Initialize counters: \(\text{num\_violations} \leftarrow 0\), \(\text{total\_triples} \leftarrow 0\), and empty list \(\text{violations} \leftarrow []\)

\STATE Randomly sample \(S\) unique triples \(\{(i, j, k)\}\) from \(\{1, \ldots, n\}\) without replacement

\FOR{each sampled triple \((i, j, k)\)}
    \STATE Extract distances \(d_{ij} = D[i,j]\), \(d_{jk} = D[j,k]\), and \(d_{ik} = D[i,k]\)
    \STATE Compute violation terms:
    \[
    v_1 = d_{ik} - \max(d_{ij}, d_{jk})
    \]
    \[
    v_2 = d_{ij} - \max(d_{ik}, d_{jk})
    \]
    \[
    v_3 = d_{jk} - \max(d_{ij}, d_{ik})
    \]
    \STATE Compute violation \(\nu = \max(v_1, v_2, v_3, 0)\)
    \IF{\(\nu > \epsilon\)}
        \STATE Append \(\nu\) to \(\text{violations}\)
        \STATE Increment \(\text{num\_violations} \leftarrow \text{num\_violations} + 1\)
    \ENDIF
    \STATE Increment \(\text{total\_triples} \leftarrow \text{total\_triples} + 1\)
\ENDFOR

\text{// Compute statistics over violations:}

\STATE $\text{max\_violation} = \max(\text{violations}) \text{ if } \text{violations} \neq \varnothing \text{ else } 0$

\STATE $\text{avg\_violation} = \text{mean}(\text{violations}) \text{ if } \text{violations} \neq \varnothing \text{ else } 0$

\STATE $\text{std\_violation} = \text{std}(\text{violations}) \text{ if } \text{violations} \neq \varnothing \text{ else } 0$


\RETURN \text{max\_violation}, \text{avg\_violation}, \text{std\_violation}

\end{algorithmic}
\end{algorithm}



% We evaluate random triplets of points and determine whether their pairwise distances satisfy the ultrametric inequality, a stricter form of the triangle inequality that enforces a strict hierarchical structure, to determine whether the distance matrix of an embedding is ultrametric.


\subsection{Neighbor Joining}

We use another metric called Neighbor Joining (NJ) to gauge how closely the pairwise distance structure of the embeddings follows a tree-like geometry. We employed a Neighbor Joining (NJ)-based analysis produced from the NJ Q-matrix~\citeauthor{trees1987neighbor}. 

Given an $n \times n$ distance matrix $D$, the NJ algorithm defines the Q-matrix as
\begin{equation}
\scriptsize
Q(i,j) = (n-2)D(i,j) - \sum_{k=1}^{n} D(i,k) - \sum_{k=1}^{n} D(j,k)
\end{equation}
where pairs with the most negative $Q(i,j)$ values are preferentially joined during tree construction. 
Rather than explicitly reconstructing a phylogenetic tree, 
we use the NJ Q-matrix, which transforms the distance matrix into a form that highlights how strongly each pair of points supports hierarchical merging. Pairs with highly negative Q-values are those most consistent with additive tree geometry. This makes the Q-matrix a useful diagnostic tool for assessing tree-likeness directly from distances.

Algorithm~\ref{alg_nj_scores} shows the computation that uses the Neighbor-Joining (NJ) Q-matrix to compute tree-likeness statistics from a given distance matrix. By modifying pairwise distances based on their global relationships, it first creates the Q-matrix. After that, all distinct pairwise Q-values are extracted, converted to absolute values, and their distribution is summarized using the maximum, mean, and standard deviation. The degree to which the data exhibit a hierarchical or tree-like structure is concisely captured by these statistics.

To obtain a global measure of hierarchical consistency, we summarize the distribution of all off-diagonal Q-values by reporting their maximum, mean, and standard deviation after normalization. The lower mean and maximum values indicate that the distance matrix more closely follows the NJ criterion and better approximates a tree metric. Combined with other geometric indicators such as hyperbolicity and curvature, these scores provide a robust measure of how strongly an embedding reflects hierarchical structure. Rather than constructing the full tree, we summarize the distribution of Q-matrix values to capture the embedding's overall consistency with a tree-structured geometry.
% % \textcolor{red}{
% Neighbor Joining leverages this structure algorithmically: when distances approximate a tree metric, the NJ Q-matrix exhibits strongly negative minima, allowing reliable reconstruction of hierarchical relationships. Consequently, low $\delta$-hyperbolicity, negative Ricci curvature, and favorable NJ scores represent mutually reinforcing signatures of tree-like geometry.
% % }

\begin{algorithm}[h]
\caption{Neighbor Joining–Based Tree-Likeness Scoring}
\label{alg_nj_scores}
\noindent
% \scriptsize
\small
\begin{tabular}{p{1cm}p{7cm}}
\textbf{Input:}  & Distance matrix \(D \in \mathbb{R}^{n \times n}\) \\
\textbf{Output:} & Dictionary containing NJ score statistics
\end{tabular}

\begin{algorithmic}[1]{}

\STATE Let \(n \leftarrow\) number of samples in \(D\)

\IF{\(n < 3\)}
    \RETURN \(\text{nj\_max} = 0,\ \text{nj\_avg} = 0,\ \text{nj\_std} = 0\)
\ENDIF

\STATE Compute the Neighbor Joining Q-matrix \(Q \in \mathbb{R}^{n \times n}\):
\[
Q(i,j) = (n-2)D(i,j) - \sum_{k=1}^{n} D(i,k) - \sum_{k=1}^{n} D(j,k)
\]

\STATE Initialize empty list \(\mathcal{Q} \leftarrow [ \ ]\)

\FOR{\(i = 1\) to \(n\)}
    \FOR{\(j = i+1\) to \(n\)}
        \STATE Append \(Q(i,j)\) to \(\mathcal{Q}\)
    \ENDFOR
\ENDFOR

\STATE Normalize Q-values: \(\mathcal{Q} \leftarrow \{\, |q| \mid q \in \mathcal{Q} \,\}\)

\STATE \(\text{nj\_max} \leftarrow \max(\mathcal{Q})\)

\STATE \(\text{nj\_avg} \leftarrow \text{mean}(\mathcal{Q})\)

\STATE \(\text{nj\_std} \leftarrow \text{std}(\mathcal{Q})\)

\RETURN \(\text{nj\_max}, \text{nj\_avg}, \text{nj\_std}\)

\end{algorithmic}
\end{algorithm}


\subsection{Ollivier-Ricci Curvature}
To capture local geometric properties complementary to global hyperbolicity, we computed the discrete Ollivier-Ricci curvature~\citeauthor{ollivier2009ricci} on the k-nearest neighbor graph constructed from embeddings (see Algorithm~\ref{alg_Ricci}). 
For adjacent points $x_i$ and $x_j$, the Ollivier-Ricci curvature is defined as:
\begin{equation}
\scriptsize
\kappa(x_i, x_j) = 1 - \frac{W_1(\mu_i, \mu_j)}{d(x_i, x_j)}
\end{equation}
where $W_1(\mu_i, \mu_j)$ is the Wasserstein-1 distance between probability distributions $\mu_i$ and $\mu_j$ centered at $x_i$ and $x_j$ respectively, typically defined as uniform distributions over their k-nearest neighbors~\citeauthor{ni2015ricci}. Negative curvature ($\kappa < 0$) indicates the local exponential divergence characteristic of hyperbolic geometry and tree-like structure~\citeauthor{sandhu2015graph}, while positive curvature suggests clustering or community structure. We report the mean, standard deviation, minimum, and proportion of edges with negative curvature in the
k-NN graph.

% \begin{algorithm}[h!]
% \caption{Compute Ollivier-Ricci Curvature}
% \label{alg_Ricci}
% % \begin{algorithmic}[1]
% % \REQUIRE Embedding matrix $X \in \mathbb{R}^{n \times d}$, number of neighbors $k$
% % \ENSURE Mean Ricci curvature $\bar{\kappa}$, standard deviation $\sigma_\kappa$, minimum $\kappa_{min}$, negative ratio $r_{neg}$
% \noindent
% \begin{tabular}{p{1cm}p{7cm}}
% \textbf{Input:}  & Embedding matrix $X \in \mathbb{R}^{n \times d}$, number of neighbors $k$ \\
% \textbf{Output:} & Mean Ricci curvature $\bar{\kappa}$, standard deviation $\sigma_\kappa$, 
%         minimum $\kappa_{min}$, negative ratio $r_{neg}$
% \end{tabular}
% \begin{algorithmic}[1]
% \STATE Construct k-NN graph $G = (V, E)$ from $X$
% \STATE Initialize $\kappa\_values \leftarrow []$
% \FOR{each edge $(i,j) \in E$}
%     \STATE $N_i \leftarrow$ k-nearest neighbors of $x_i$
%     \STATE $N_j \leftarrow$ k-nearest neighbors of $x_j$
%     \STATE Define uniform distributions: $\mu_i = \text{Uniform}(N_i)$, 
%            $\mu_j = \text{Uniform}(N_j)$
%     \STATE Compute Wasserstein-1 distance: $W_1 = W_1(\mu_i, \mu_j)$
%     \STATE $\kappa(i,j) \leftarrow 1 - \frac{W_1}{d(x_i, x_j)}$
%     \STATE Append $\kappa(i,j)$ to $\kappa\_values$
% \ENDFOR
% \STATE $\bar{\kappa} \leftarrow \text{mean}(\kappa\_values)$
% \STATE $\sigma_\kappa \leftarrow \text{std}(\kappa\_values)$
% \STATE $\kappa_{min} \leftarrow \min(\kappa\_values)$
% \STATE $r_{neg} \leftarrow \frac{\#\{\kappa < 0\}}{|\kappa\_values|}$
% \RETURN $\bar{\kappa}, \sigma_\kappa, \kappa_{min}, r_{neg}$
% \end{algorithmic}
% \end{algorithm}


\begin{algorithm}[h!]
\caption{Compute Ollivier-Ricci Curvature}
\label{alg_Ricci}
\small
% \begin{algorithmic}[1]
% \REQUIRE Embedding matrix $X \in \mathbb{R}^{n \times d}$, number of neighbors $k$
% \ENSURE Mean Ricci curvature $\bar{\kappa}$, standard deviation $\sigma_\kappa$, minimum $\kappa_{min}$, negative ratio $r_{neg}$
\noindent
\begin{tabular}{p{1cm}p{7cm}}
\textbf{Input:}  & Embedding matrix $X \in \mathbb{R}^{n \times d}$, number of neighbors $k$ \\
\textbf{Output:} & Mean Ricci curvature $\bar{\kappa}$, standard deviation $\sigma_\kappa$, 
        minimum $\kappa_{min}$, negative ratio $r_{neg}$
\end{tabular}
\begin{algorithmic}[1]
\STATE Construct k-NN graph $G = (V, E)$ from $X$
\STATE Initialize $\kappa\_values \leftarrow []$
\FOR{each edge $(i,j) \in E$}
    \STATE $N_i \leftarrow$ k-nearest neighbors of $x_i$
    \STATE $N_j \leftarrow$ k-nearest neighbors of $x_j$
    \STATE Define uniform distributions: $\mu_i = \text{Uniform}(N_i)$, 
           $\mu_j = \text{Uniform}(N_j)$
    \STATE Compute Wasserstein-1 distance: $W_1 = W_1(\mu_i, \mu_j)$
    \STATE $\kappa(i,j) \leftarrow 1 - \frac{W_1}{d(x_i, x_j)}$
    \STATE Append $\kappa(i,j)$ to $\kappa\_values$
\ENDFOR
\STATE $\bar{\kappa} \leftarrow \text{mean}(\kappa\_values)$
\STATE $\sigma_\kappa \leftarrow \text{std}(\kappa\_values)$
\STATE $\kappa_{min} \leftarrow \min(\kappa\_values)$
\STATE $r_{neg} \leftarrow \frac{\#\{\kappa < 0\}}{|\kappa\_values|}$
\RETURN $\bar{\kappa}, \sigma_\kappa, \kappa_{min}, r_{neg}$
\end{algorithmic}
\end{algorithm}


\subsection{Persistent homology}

We employ 0-dimensional persistent homology~\citeauthor{edelsbrunner2008persistent, carlsson2009topology} to quantify the topological structure of embeddings across scales (see Algorithm~\ref{alg_Persistent}). Persistent homology tracks connected components as we increase the distance threshold $\epsilon$ from 0 to $\infty$. Each connected component has a "birth" time (when it first appears) and ``death" time (when it merges with another component). The persistence lifetime $\ell = \text{death} - \text{birth}$ measures the stability of that topological feature~\citeauthor{chazal2016structure}.
\\
For tree-like structures with clear hierarchical organization, components merge rapidly, yielding short persistence lifetimes. In contrast, flat or diffuse geometries exhibit long persistence lifetimes as components remain separated across many scales. We compute the persistence diagram using the Ripser algorithm~\citeauthor{bauer2021ripser} and report:
\begin{itemize}
\item Mean persistence lifetime: $\bar{\ell} = \frac{1}{n-1}\sum_{i=1}^{n-1} \ell_i$
\item Maximum persistence lifetime: $\max_i \ell_i$
\item Persistence entropy: $H = -\sum_{i=1}^{n-1} p_i \log p_i$ where 
      $p_i = \ell_i / \sum_j \ell_j$
\end{itemize}
Lower persistence values indicate stronger hierarchical consolidation.

\begin{algorithm}[h!]
\caption{Compute Persistent Homology Features}
\label{alg_Persistent}
\small
\noindent
\begin{tabular}{p{1cm}p{7cm}}
\textbf{Input:}  & Distance matrix $D \in \mathbb{R}^{n \times n}$ \\
\textbf{Output:} & Mean persistence $\bar{\ell}$, max persistence $\ell_{max}$, persistence entropy $H$
\end{tabular}
\begin{algorithmic}[1]
% \REQUIRE Distance matrix $D \in \mathbb{R}^{n \times n}$
% \ENSURE Mean persistence $\bar{\ell}$, max persistence $\ell_{max}$, persistence entropy $H$
\STATE Compute 0-dimensional persistence diagram using Ripser: 
       $\{(\text{birth}_i, \text{death}_i)\}_{i=1}^{n-1}$
\STATE Compute lifetimes: $\ell_i = \text{death}_i - \text{birth}_i$
\STATE $\bar{\ell} \leftarrow \frac{1}{n-1}\sum_{i=1}^{n-1} \ell_i$
\STATE $\ell_{max} \leftarrow \max_{i} \ell_i$
\STATE Normalize: $p_i \leftarrow \frac{\ell_i}{\sum_{j=1}^{n-1} \ell_j}$
\STATE $H \leftarrow -\sum_{i=1}^{n-1} p_i \log p_i$
\RETURN $\bar{\ell}, \ell_{max}, H$
\end{algorithmic}
\end{algorithm}



\subsection{Fractal Dimension}

We estimate intrinsic dimensionality using the correlation dimension~\citeauthor{grassberger1983measuring}, derived from the scaling behavior of the correlation integral for points ${x_i}$. Details are mentioned in Algorithm~\ref{alg_Fractal}.  Correlation integral for points ${x_i}$ can be defined by:
\begin{equation}
\scriptsize
C(r) = \lim_{n \to \infty} \frac{1}{n^2} \sum_{i \neq j} \mathbb{I}(d(x_i, x_j) < r)
\end{equation}
In the scaling region, $C(r) \propto r^{d_f}$, so $d_f$ can be estimated from the slope of $\log C(r)$ vs $\log r$~\citeauthor{camastra2016intrinsic}. We use the two-nearest-neighbor method~\citeauthor{facco2017estimating} for robust estimation:
\begin{equation}
\scriptsize
d_f = \left[ \frac{1}{n} \sum_{i=1}^n \log\left(\frac{r_2(x_i)}{r_1(x_i)}\right) \right]^{-1}
\end{equation}
where $r_1(x_i)$ and $r_2(x_i)$ are the distances to the first and second nearest neighbors of $x_i$. Lower fractal dimension signals embeddings concentration near low-dimensional manifolds, whereas higher values indicate a more diffuse, isotropic spread in high-dimensional space~\citeauthor{levina2004maximum}.


\begin{algorithm}[h!]
\caption{Estimate Fractal Dimension}
\label{alg_Fractal}
\small
\noindent
\begin{tabular}{p{1cm}p{7cm}}
\textbf{Input:}  & Embedding matrix $X \in \mathbb{R}^{n \times d}$ \\
\textbf{Output:} & Mean fractal dimension $\bar{d}_f$, standard deviation $\sigma_{d_f}$, median $\tilde{d}_f$
\end{tabular}
\begin{algorithmic}[1]
% \REQUIRE Embedding matrix $X \in \mathbb{R}^{n \times d}$
% \ENSURE Mean fractal dimension $\bar{d}_f$, standard deviation $\sigma_{d_f}$, median $\tilde{d}_f$
\STATE $d_f\_values \leftarrow []$
\FOR{each point $x_i \in X$}
    \STATE Find distances to 2 nearest neighbors: $r_1(x_i), r_2(x_i)$
    \STATE $d_{f,i} \leftarrow \left[\log\left(\frac{r_2(x_i)}{r_1(x_i)}\right)\right]^{-1}$
    \STATE Append $d_{f,i}$ to $d_f\_values$
\ENDFOR
\STATE $\bar{d}_f \leftarrow \text{mean}(d_f\_values)$
\STATE $\sigma_{d_f} \leftarrow \text{std}(d_f\_values)$
\STATE $\tilde{d}_f \leftarrow \text{median}(d_f\_values)$
\RETURN $\bar{d}_f, \sigma_{d_f}, \tilde{d}_f$
\end{algorithmic}
\end{algorithm}

\section{Experimental Setup}
\label{sec_es}
We performed the experiments on a computer running 64-bit Windows 10 with an Intel(R) Core i5 processor at 2.10 GHz and 32 GB of RAM. For experiments, we randomly split the data into 60-10-30\% for training, validation, and testing. We repeat the experiments $5$ times and report the average and standard deviation.
For the sake of reproducibility, we have made code available online~\footnote{\url{https://github.com/pchourasia1/Hyperbolicity_LLM_Embeddings}}.
% Our Data and code are available online for reproducibility~\footnote{available in the published version}.

\subsection{Data Statistics}
\label{sub_sec_data_stats}
We use $3$ real-world protein sequence datasets. The summary of the datasets used for the experimentation is given below:

\subsubsection{PDB186}
We use a benchmark dataset from the literature (PDB186) for predicting DNA-protein binding, a binary classification task~\citeauthor{lou2014sequence}. More formally, the dataset contains DNA-binding and non-binding protein sequences. The research task is to predict whether the DNA and protein sequences will bind. The minimum, maximum, and average sequence length in the data are 64, 1323, and 264.693, respectively. The total number of sequences is 186, of which 93 bind and 93 do not.

% The binding prediction problem in PDB186 is useful for health recommendation systems in the following ways:
% \paragraph{\textbf{Improving Predictive Models for Personalized Medicine:}}
% DNA-Protein interactions are fundamental in understanding various biological processes, including gene expression and regulation. By predicting whether a DNA sequence will bind to a protein, the algorithm can help identify key genetic markers or mutations that might influence treatment effectiveness or susceptibility to certain diseases. This information can be used to develop more personalized and precise health recommendations, such as tailoring medication or treatment plans based on an individual’s genetic profile.

% \paragraph{\textbf{Drug Discovery and Development:}} Predicting DNA-Protein interactions is critical in drug discovery, where identifying target proteins that bind to specific DNA sequences can lead to the development of new therapies. Health recommendation systems can incorporate these findings to suggest drugs that are more likely to be effective for specific genetic conditions, thereby improving treatment outcomes.

% \paragraph{\textbf{Genetic Counseling and Risk Assessment:}} By accurately predicting DNA-Protein binding, the algorithm can contribute to risk assessment models in genetic counseling. For example, individuals with a genetic predisposition to certain diseases could receive health recommendations that include lifestyle changes, preventive measures, or early interventions based on their unique DNA-Protein interaction profiles.

% \paragraph{\textbf{Support for Clinical Decision-Making:}} In clinical settings, the algorithm's predictions can assist healthcare professionals in making more informed decisions. For instance, understanding how a patient’s genetic makeup influences DNA-Protein binding can help clinicians predict the patient’s response to certain therapies and adjust treatment plans accordingly.


\subsubsection{Breast Cancer}
Our Membranolytic anticancer peptides (ACPs) dataset~\citeauthor{Grisoni2019} contains 949 protein sequences annotated for anticancer activity against breast cancer cell lines. The target labels are categorized into ``very active'', ``moderately active'', ``experimental inactive'', and ``virtual inactive'' groups. Table~\ref{tab_data_Stats_BC} shows the distribution in detail.


\begin{table}[!t]
      \centering
      \resizebox{0.7\textwidth}{!}{
         \begin{tabular}{lcccc}
    \toprule
    ACPs Category & Count & Min. & Max. & Average \\
    \midrule \midrule
        Inactive-Virtual & 750 & 8 & 30 & 16.64 \\
        Moderate Active & 98 & 10 & 38 & 18.44 \\
        Inactive-Experimental & 83 & 5 & 38 & 15.02 \\
        Very Active & 18 & 13 & 28 & 19.33 \\
        \midrule \midrule
        Total  & 949 & - & - & -\\
        \bottomrule
    \end{tabular}
    }
    % \caption*{Breast cancer data}
    % \label{tab_data_Stats_BC}
    \caption{Dataset statistics for the Breast cancer data. Columns represent the minimum, maximum, and average sequence lengths.}
    \label{tab_data_Stats_BC}
\end{table}


% \begin{table}[h!]
%     \centering
%     \resizebox{0.49\textwidth}{!}{
%     \begin{tabular}{@{\extracolsep{6pt}}lccccccc}
%     \toprule
%      & & \multicolumn{3}{c}{Peptide Sequence Length} & \multicolumn{3}{c}{Number of Sequences} \\
%      \cmidrule{3-5}  \cmidrule{6-8}
%     ACPs Category & Count & Min. & Max. & Average & Training & Validation & Testing \\
%     \midrule \midrule
%         Inactive-Virtual & 750 & 8 & 30 & 16.64 & 540 & 60 & 150\\
%         Moderate Active & 98 & 10 & 38 & 18.44 & 71 & 7 & 19 \\
%         Inactive-Experimental & 83 & 5 & 38 & 15.02 & 61 & 6 & 16 \\
%         Very Active & 18 & 13 & 28 & 19.33 & 14 & 1 & 3 \\
%         \midrule \midrule
%         Total  & 949 & - & - & - & - & - & - \\
%         \bottomrule
%     \end{tabular}
%     }
%     \caption{ACPs dataset distribution based on their respective activity on the breast cancer cell line. The minimum, maximum, and average sequence lengths for each category are also reported, along with the counts of sequences in the test, validation, and train sets for each category. }
%     \label{tab_data_Stats}
% \end{table}

% The target labels are categorized into ``very active'' (750 sequences), ``moderately active'' (98 sequences), ``experimental inactive'' (83 sequences), and ``virtual inactive'' (18 sequences) groups. Most sequences range from about 5 to 38 amino acids in length. This dataset provides a diverse set of short peptide sequences spanning a broad range of activity levels. Table~\ref{tab_data_Stats_BC} (in the appendix) shows the distribution in detail.

% Our Membranolytic anticancer peptides (ACPs) dataset~\citeauthor{Grisoni2019} contains information about the peptides (protein sequences) and their anticancer activity (target labels) on breast cancer cell lines. The target labels are categorized into ``very active'', ``moderately active'', ``experimentally inactive'', and ``virtual inactive'' groups. This dataset contains $949$ peptide sequences distributed among the four categories, as shown in Table~\ref{tab_data_Stats_BC}.


% \begin{table}[!t]
%     \caption{Dataset statistics for the Breast cancer data. Columns represent the min., max., and average lengths of the sequence.}
%     \label{tab_data_Stats_BC}
%       \centering
%       \resizebox{0.35\textwidth}{!}{
%          \begin{tabular}{lcccc}
%     \toprule
%     ACPs Category & Count & Min. & Max. & Average \\
%     \midrule \midrule
%         Inactive-Virtual & 750 & 8 & 30 & 16.64 \\
%         Moderate Active & 98 & 10 & 38 & 18.44 \\
%         Inactive-Experimental & 83 & 5 & 38 & 15.02 \\
%         Very Active & 18 & 13 & 28 & 19.33 \\
%         \midrule \midrule
%         Total  & 949 & - & - & -\\
%         \bottomrule
%     \end{tabular}
%     }
%     % \caption*{Breast cancer data}
%     % \label{tab_data_Stats_BC}
% \end{table}


% \begin{table}[h!]
%     \centering
%     \resizebox{0.9\textwidth}{!}{
%     \begin{tabular}{@{\extracolsep{6pt}}lccccccc}
%     \toprule
%      & & \multicolumn{3}{c}{Peptide Sequence Length} & \multicolumn{3}{c}{Number of Sequences} \\
%      \cmidrule{3-5}  \cmidrule{6-8}
%     ACPs Category & Count & Min. & Max. & Average & Training & Validation & Testing \\
%     \midrule \midrule
%         Inactive-Virtual & 750 & 8 & 30 & 16.64 & 540 & 60 & 150\\
%         Moderate Active & 98 & 10 & 38 & 18.44 & 71 & 7 & 19 \\
%         Inactive-Experimental & 83 & 5 & 38 & 15.02 & 61 & 6 & 16 \\
%         Very Active & 18 & 13 & 28 & 19.33 & 14 & 1 & 3 \\
%         \midrule \midrule
%         Total  & 949 & - & - & - & - & - & - \\
%         \bottomrule
%     \end{tabular}
%     }
%     \caption{ACPs dataset distribution based on their respective activity on the breast cancer cell line. The minimum, maximum, and average sequence lengths for each category are also reported, along with the counts of sequences in the test, validation, and train sets for each category. }
%     \label{tab_data_Stats}
% \end{table}



\subsubsection{Lung Cancer}
The Membranolytic anticancer peptides (ACPs) dataset~\citeauthor{Grisoni2019} contains 901 sequences annotated for anticancer activity against Lung cancer cell lines.  The target labels are classified into four groups: ``very active," ``moderately active," ``experimental inactive," and ``virtual inactive" groups. Table~\ref{tab_data_Stats_Lung_C} shows the distribution in detail.



\begin{table}[!t]
      \centering
      \resizebox{0.7\textwidth}{!}{
         \begin{tabular}{lcccc}
    \toprule
    ACPs Category & Count & Min. & Max. & Average \\
    \midrule \midrule
        Inactive-Virtual & 750 & 8 & 30 & 16.64 \\
        Moderate Active & 75 & 11 & 38 & 17.76 \\
        Inactive-Experimental & 52 & 5 & 38 & 14.5 \\
        Very Active & 24 & 13 & 28 & 20.70 \\
        \midrule \midrule
        Total  & 901 & - & - & -\\
        \bottomrule
    \end{tabular}
    }
     \caption{Dataset statistics for the Lung cancer data. Columns represent the min., max., and average lengths of the sequence.}
     \label{tab_data_Stats_Lung_C}
\end{table}


% The target labels are categorized into ``very active'' (750 sequences), ``moderately active'' (75 sequences), ``experimental inactive'' (52 sequences), and ``virtual inactive'' (24 sequences) groups.
% This collection provides a broad spectrum of short peptide sequences with varying levels of anticancer activity. Table~\ref{tab_data_Stats_LC_BC} (in the appendix) shows the distribution in detail.

% The dataset on Membranolytic anticancer peptides (ACPs)~\citeauthor{Grisoni2019} provides details regarding peptide sequences and their corresponding anticancer effectiveness against breast and lung cancer cell lines. The target labels are classified into four groups: ``very active," ``moderately active," ``experimental inactive," and ``virtual inactive." In total, the dataset comprises 949 and 901 peptide sequences for breast and lung cancer, respectively. Table~\ref{tab_data_Stats_LC_BC} shows the distribution.




% \begin{table}[!t]
%  \caption{Dataset statistics for the Lung cancer data. Columns represent the min., max., and average lengths of the sequence.}
%  \label{tab_data_Stats_LC}
%       \centering
%       \resizebox{0.35\textwidth}{!}{
%          \begin{tabular}{lcccc}
%     \toprule
%     ACPs Category & Count & Min. & Max. & Average \\
%     \midrule \midrule
%         Inactive-Virtual & 750 & 8 & 30 & 16.64 \\
%         Moderate Active & 75 & 11 & 38 & 17.76 \\
%         Inactive-Experimental & 52 & 5 & 38 & 14.5 \\
%         Very Active & 24 & 13 & 28 & 20.70 \\
%         \midrule \midrule
%         Total  & 901 & - & - & -\\
%         \bottomrule
%     \end{tabular}
%     }
%     % \caption*{Lung cancer data}
%     % \label{tab_data_Stats_LC}
%     % \caption{Dataset statistics for the Breast cancer and  Lung cancer data. Columns represent the min., max., and average lengths of the peptide sequence.}
%     \label{tab_data_Stats_LC_BC}
% \end{table}

% \subsection{Evaluation}
% To evaluate the results, we analyze the $\delta$-hyperbolicity, ultrametricity, and Neighbor Joining for LLM embeddings. To understand the utility of these metrics on the different datasets and to understand the structural properties. 
% \textcolor{red}{we generate synthetic data in sphere space, dense graph, and Poincaré space.} 
% Moreover, we perform classification on the TAPE embeddings using classical machine learning classifiers like support vector machine (SVM), Decision Tree (DT), Random Forest (RF), Logistic Regression (LR), Multi-Layer Perceptron (MLP), K-nearest neighbors (KNN), and Naive Bayes (NB). We divide the data into 70-30 \% training and testing, compute the result $5$ times, and report average and standard deviation values for accuracy, precision, recall, weighted F1 score, macro F1 score, ROC-AUC, and classifier training runtime.

\subsection{Classification}
For the classification task, we use classifiers like SVM, Naive Bayes, MLP, KNN, Random Forest, Logistic Regression, and Decision Tree. Evaluation metrics include accuracy, precision, recall, weighted F1, macro F1, ROC-AUC, and training runtime. 
% The preprocessed data and code will be available (in the published version).

\subsection{Clustering}
For the clustering task, we use $k$-means~\citeauthor{kmeans}, $k$-modes~\citeauthor{huang-1998-extensions}, and the agglomerative clustering algorithm~\citeauthor{ward1963hierarchical}. The evaluation metrics used were the Silhouette Coefficient~\citeauthor{rousseeuw1987silhouettes}, Calinski-Harabasz Score~\citeauthor{calinski1974dendrite}, and the Davies-Bouldin Score~\citeauthor{davies1979cluster}.

\subsubsection{Algorithms}

\paragraph{$k$-means~\citeauthor{kmeans}}
The $k$-means~\citeauthor{kmeans} clusters
objects based on the Euclidean mean. The $k$-means clustering partitions the data into $K$ groups by repeatedly assigning points to the nearest centroid and updating the centroids as the mean of their assigned points. This continues until the assignments stabilize, resulting in compact, well-separated clusters~\citeauthor{kanungo2002efficient}.

\paragraph{$k$-modes~\citeauthor{huang-1998-extensions}}
We also perform clustering with $k$-modes~\citeauthor{huang-1998-extensions}, which is a variant of $k$-means
using \textit{modes} instead of \textit{means}.  Here, pairs of
objects are subject to a dissimilarity measure (\eg, Hamming distance) rather than to the Euclidean mean.

\paragraph{Agglomerative~\citeauthor{ward1963hierarchical}}
All agglomerative hierarchical clustering methods start from a proximity matrix of either similarities or distances between elements ~\citeauthor{nei1978estimation}. Agglomerative clustering is a bottom-up algorithm that starts with each data point as its own cluster. It repeatedly merges the closest pairs of clusters based on a chosen distance metric until all points are grouped into a single cluster or a set number of clusters~\citeauthor{mojena1977hierarchical}. 


\subsubsection{Evaluation Metric}

\paragraph{Silhouette Coefficient~\citeauthor{rousseeuw1987silhouettes}}
Measures how similar a point is to its own cluster versus similarity with other clusters. Scores range from $-1$ to $1$, with higher values indicating better clustering.

\paragraph{Calinski–Harabasz Score~\citeauthor{calinski1974dendrite}}
Evaluates clustering quality using between-cluster and within-cluster dispersion. Higher values indicate more distinct, well-separated clusters.

\paragraph{Davies–Bouldin Score~\citeauthor{davies1979cluster}}
Computes the ratio of within-cluster to between-cluster distances. Lower values indicate better separation between clusters.



\subsection{Visualization}
A popular visualization technique, named t-SNE~\citeauthor{van2008visualizing}, is employed to visualize the feature vectors generated by the SeqVec~\citeauthor{heinzinger2019modeling} method. 
% t-SNE has been widely used to visualize biological data. 
The t-SNE plots are shown in Figure~\ref{tsne_pdb186}  for PDB186, Figure~\ref{tsne_lung_cancer} for Lung Cancer, and Figure~\ref{tsne_breast_cancer} for breast cancer datasets. The main idea for reporting the t-SNE plots is to assess whether there is a clear class separation between the labels in the datasets using embeddings generated by the respective LLMs. 
We observe that clusters overlap for Seqvec in the Lung cancer and breast cancer datasets, indicating that no clear decision boundary exists in the data.




\begin{figure*}[h!]
     \centering
    \begin{subfigure}[b]{0.23\textwidth}
         \centering
         \includegraphics[scale=0.075]{Figures/tsne/PDB_186/seqvec_padded_embeddings_PDB_186.png}
         \caption{Seqvec}
         \label{tsne_seqvec_pdb186}
     \end{subfigure}
     \hfill
     \begin{subfigure}[b]{0.23\textwidth}
         \centering
         \includegraphics[scale=0.075]{Figures/tsne/PDB_186/ESM2_PDB186_Embeddings.png}
         \caption{ESM2}
         \label{tsne_esm2_pdb186}
     \end{subfigure}
     % \hfill
     % \\
     \begin{subfigure}[b]{0.23\textwidth}
         \centering
         \includegraphics[scale=0.075]{Figures/tsne/PDB_186/PDB186_TAPE_Embedding.png}
         \caption{TAPE}
         \label{tsne_tape_pdb186}
     \end{subfigure}
          \hfill
     \begin{subfigure}[b]{0.23\textwidth}
         \centering
         \includegraphics[scale=0.075]{Figures/tsne/PDB_186/ProtT5_embeddings_PDB186.png}
         \caption{ProtT5}
         \label{tsne_protT5_pdb186}
     \end{subfigure}
     \includegraphics[width=0.65\textwidth]{Figures/tsne/PDB_186/PDB186_legends.png}
     \caption{t-SNE visualization of different embeddings from\textbf{ PDB 186 dataset}. The Figure is best seen in color.}
     \label{tsne_pdb186}
\end{figure*}

\begin{figure*}[h!]
     \centering
    \begin{subfigure}[b]{0.23\textwidth}
         \centering
         \includegraphics[scale=0.075]{Figures/tsne/LC/SeqVec_Lungs_Cancer_901.png}
         \caption{Seqvec}
         \label{tsne_seqvec_LC}
     \end{subfigure}
     \hfill
     \begin{subfigure}[b]{0.23\textwidth}
         \centering
         \includegraphics[scale=0.075]{Figures/tsne/LC/ESM2_Lungs_Cancer_901.png}
         \caption{ESM2}
         \label{tsne_esm2_LC}
     \end{subfigure}
     % \hfill
     % \\
     \begin{subfigure}[b]{0.23\textwidth}
         \centering
         \includegraphics[scale=0.075]{Figures/tsne/LC/TAPE_Lungs_Cancer_901.png}
         \caption{TAPE}
         \label{tsne_tape_LC}
     \end{subfigure}
          \hfill
     \begin{subfigure}[b]{0.23\textwidth}
         \centering
         \includegraphics[scale=0.075]{Figures/tsne/LC/ProtT5_Lungs_Cancer_901.png}
         \caption{ProtT5}
         \label{tsne_protT5_LC}
     \end{subfigure}
     \includegraphics[width=0.75\textwidth]{Figures/tsne/LC/Lung_Cancer_legends.png}
     \caption{t-SNE visualization of different embeddings from\textbf{ Lung Cancer dataset}. The Figure is best seen in color.}
     \label{tsne_lung_cancer}
\end{figure*}

\begin{figure*}[h!]
     \centering
    \begin{subfigure}[b]{0.23\textwidth}
         \centering
         \includegraphics[scale=0.075]{Figures/tsne/BC/SeqVec_Breast_Cancer_949.png}
         \caption{Seqvec}
         \label{tsne_seqvec_BC}
     \end{subfigure}
     \hfill
     \begin{subfigure}[b]{0.23\textwidth}
         \centering
         \includegraphics[scale=0.075]{Figures/tsne/BC/ESM2_Breast_Cancer_949.png}
         \caption{ESM2}
         \label{tsne_esm2_BC}
     \end{subfigure}
     % \hfill
     % \\
     \begin{subfigure}[b]{0.23\textwidth}
         \centering
         \includegraphics[scale=0.075]{Figures/tsne/BC/TAPE_Breast_Cancer_949.png}
         \caption{TAPE}
         \label{tsne_tape_BC}
     \end{subfigure}
          \hfill
     \begin{subfigure}[b]{0.23\textwidth}
         \centering
         \includegraphics[scale=0.075]{Figures/tsne/BC/ProtT5_Breast_Cancer_949.png}
         \caption{ProtT5}
         \label{tsne_protT5_BC}
     \end{subfigure}
     \includegraphics[width=0.75\textwidth]{Figures/tsne/BC/BC_legends.png}
     \caption{t-SNE visualization of different embeddings from\textbf{ Breast Cancer dataset}. The Figure is best seen in color.}
     \label{tsne_breast_cancer}
\end{figure*}

\section{Results And Discussion}
\label{sec_results}

In the first step, we simulate the data to understand the hyperbolicity values for known structures using the following synthetic data.
\paragraph{Sphere Space Data Computation} 
% We sample random points on a high-dimensional unit sphere by normalizing Gaussian vectors and computing Euclidean distances. Hyperbolicity is estimated from the resulting distance matrix.
To simulate data in the sphere space, we generate $n$ random points on the surface of a unit sphere in a high-dimensional space. Each point is sampled from a standard normal distribution in $\mathbb{R}^d$, and subsequently normalized to lie on the unit sphere by dividing each point by its Euclidean norm. The distance matrix for these points is computed using the Euclidean distance metric. The $\delta$-hyperbolicity values for this space are then calculated by evaluating the shortest path distances between all possible pairs of points, followed by applying the $\delta$-hyperbolicity calculation method described in Section~\ref{sec_PAa}.

\paragraph{Dense Graph Data Computation} 
A dense graph is generated using the Erdős-Rényi model, where $n$ nodes are connected with a high probability, $p = 0.8$. The shortest path distances are computed using the Floyd–Warshall algorithm. $\delta$-hyperbolicity is evaluated from this distance matrix, assigning a large finite value to any disconnected pairs to keep the computation valid.
% A dense graph is generated using the Erdős-Rényi model, where $n$ nodes are connected with a high probability, $p = 0.8$. The adjacency matrix of the graph is obtained, representing the presence of edges between nodes. To compute the distance matrix, the Floyd-Warshall algorithm is applied to find the shortest paths between all pairs of nodes. The $\delta$-hyperbolicity for this space is computed by analyzing the shortest path distances, with special handling for disconnected node pairs by assigning them a large finite distance to ensure valid hyperbolicity computations.

\paragraph{Poincaré Space Data Computation} 
To simulate data in the Poincaré space, we sample $n$ random points in the two-dimensional Poincaré ball by drawing radial distances $r \in [0,1)$ and angles $\theta \in [0,2\pi)$, then converting to Cartesian coordinates via $x = r\cos\theta$ and $y = r\sin\theta$. The geodesic distances of pairs are calculated using the hyperbolic distance formula, which incorporates both the Euclidean separation and the distance of each point from the origin.

% To simulate data in the Poincaré space, we generate $n$ random points within the two-dimensional Poincaré ball. Each point is characterized by a radial distance $r$ from the origin and an angular coordinate $\theta$. The radial distance $r$ is uniformly sampled from the interval $[0, 1)$, while the angular coordinate $\theta$ is uniformly sampled from the interval $[0, 2\pi)$. The Cartesian coordinates of each point are then calculated using the polar-to-Cartesian transformation: $x = r \cos(\theta)$ and $y = r \sin(\theta)$. The geodesic distance between any two points within the Poincaré ball is computed using the hyperbolic distance formula, which considers both the Euclidean distance between the points and their distances from the origin. 


% \paragraph{Hyperbolicity Computation} 
% For the sphere space, the dense graph, and the Poincaré Space, the hyperbolicity is computed using the distance matrices derived from their respective constructions. The computation involves evaluating the four-point condition for all possible quadruples of points or nodes. Specifically, for any quadruple of points, we calculate the $\delta$ value, which represents the deviation from an ideal tree metric. The average and standard deviation of these $\delta$ values are then reported in Table~\ref{tbl_synthetic} as the hyperbolicity measures for the respective metric spaces.

\begin{table}[h!]
    \centering
    \resizebox{0.6\textwidth}{!}{
    \begin{tabular}{cc}
    \toprule
         Metric Space & $\delta$ \\
         \midrule \midrule
         Sphere Space - $\delta$: & 0.32 $\pm$ 0.19 \\
        Dense Graph - $\delta$: & 0.31 $\pm$ 0.25 \\
        Poincaré Space - $\delta$: & 0.2867 $\pm$ 0.3328 \\
        \bottomrule
    \end{tabular}
    }
    \caption{Hyperbolicity $\delta$ values (average and standard deviation) for synthetically generated data.}
    \label{tbl_synthetic}
\end{table}

The differences in $\delta$-hyperbolicity across the three synthetic spaces reflect their underlying geometry as shown in Table~\ref{tbl_synthetic}. Sphere and dense graph data show moderately higher $\delta$ values of $0.32 \pm 0.19$ and $0.31 \pm 0.25$, respectively, due to curvature and high connectivity, both of which reduce tree-likeness. The Poincaré space, on the other hand, produces the lowest $\delta$ value of $0.2867 \pm 0.3328$, which is consistent with its inherently hyperbolic, tree-like structure~\citeauthor{nickel2017poincare}. This illustrates how a space's tree-like appearance is shaped by geometry.

% The variation in hyperbolicity values across different metric spaces, as shown in Table~\ref{tbl_synthetic}, can be attributed to the underlying structural properties of the data. The sphere space, with a $\delta$ value of $0.32 \pm 0.19$, exhibits slightly higher hyperbolicity, likely due to the curved nature of the space, which deviates from the tree-like structure that minimizes hyperbolicity. The dense graph, with a $\delta$ value of $0.31 \pm 0.25$, exhibits similar levels of hyperbolicity, reflecting its high connectivity and the presence of many alternative paths between nodes, which reduces its tree-likeness. In contrast, the Poincaré space, with the smallest $\delta$ value of $0.2867 \pm 0.3328$, reflects its inherent hyperbolic geometry, which naturally supports a tree-like structure with minimal deviation (this hypothesis aligns well with literature on Poincaré ball analysis~\citeauthor{nickel2017poincare}), leading to lower hyperbolicity. The differences in hyperbolicity across these spaces highlight how the geometric and topological characteristics of the data influence the extent to which the space approximates a tree metric.


% \begin{table}[h!]
%     \centering
%     \begin{tabular}{ccc}
%     \toprule
%     Dataset & $\delta_{Average}$ & $\delta_{Std.}$ \\
%     \midrule \midrule
%        PDB186  & 0.4178 & 0.8754 \\
%         Lung Cancer & 0.4973 & 1.1015 \\
%         Breast Cancer & \_ & \_ \\
%        \bottomrule
%     \end{tabular}
%     \caption{Hyperbolicity $\delta$ values (average and standard deviation) for different datasets. \textcolor{red}{SARWAN: added lung cancer, need to add more datasets and their discussion accordingly}}
%     \label{tbl_delta}
% \end{table}


% \begin{table}[h!]
%     \centering
%     \begin{tabular}{cccc}
%     \toprule
%     Embedding & $\delta_{max}$ & $\delta_{Average}$ & $\delta_{Std.}$ \\
%     \midrule \midrule
%        SeqVec & 116.7348 & 3.2018 & 6.9712 \\  
%        % seqvec is after PCA 172 components with 99 percent variance were selected
%        % ESM & 8.7726 & 0.8727 &  1.9674 \\
%        ESM2 & 6.0017 & 0.1913 & 0.4090 \\
%        TAPE & 9.6066 & 0.4178 & 0.8754 \\
%        ProtT5 & 1.0945 & 0.0418 & 0.0905 \\
%        \bottomrule
%     \end{tabular}
%     \caption{Hyperbolicity $\delta$ values (average and standard deviation) for PDB186 dataset.}
%     \label{tbl_delta}
% \end{table}

% \begin{table}[h!]
%     \centering
%     \begin{tabular}{cccc|c}
%     \toprule
%     Embedding & $max$ & $Average$ & $Std.$ & CCC \\
%     \midrule \midrule
%        Seqvec & 142.5040 &  16.6730 & 16.7692 & \_ \\
%        ESM2 & 3.5801 &  0.4888  & 0.4281 & 0.7604 \\
%        % ESM2 & 2.6213 &  0.6518  & 0.4103 \\
%        TAPE & 6.1348 &  1.0116  & 0.7845 & 0.8569 \\
%        ProtT5 & 0.9371 & 0.1301 & 0.1011 & \_ \\
       
%        \bottomrule
%     \end{tabular}
%     \caption{Ultrametricity values (max, and average) for\textbf{PDB186} datasets.}
%     \label{tbl_delta}
% \end{table}

\subsection{Hierarchical, Geometric and Topological Results}

% \textcolor{red}{Prakash: We need to explain the NJ score and discussion for Poincare distance matrix or the Euclidean Distance matrix and its impact on results. Inference of tables 4, 5, and 6}
\subsubsection{PDB186 dataset}
Table \ref{tab_PDB186_embedding_metrics} summarizes the geometric structure of the embeddings across all metrics. TAPE has the highest tree-likeness in the Poincaré space, with the lowest fractal dimension, hyperbolicity, ultrametricity, and NJ deviation. SeqVec shows the most negative Ricci curvature and the highest persistence entropy, indicating strong local hyperbolicity but weaker global hierarchy.ProtT5 exhibits significant deviations from tree structure and performs poorly in Poincaré space.
With the lowest $\delta_{\text{avg}}$, $U_{\text{avg}}$, and NJ values under Euclidean distances, ProtT5 becomes the most tree-like, and ESM performs the worst. In general, Poincaré distances yield far smaller and more stable hierarchy metrics, confirming that hyperbolic geometry captures tree structure more effectively than Euclidean space.

Overall, the results show that (i) ProtT5 works best when utilizing Euclidean distances; (ii) TAPE is the most hierarchical in hyperbolic space; and (iii) SeqVec shows great local hyperbolicity but weaker global hierarchy. Results show that the tree-likeness of the resulting representation space is significantly influenced by both the embedding model and the underlying distance geometry.

% The variation in hyperbolicity values for the PDB186 dataset, as shown in Table~\ref{tab_pdb186_distance_matrix_comparison}, can be attributed to the underlying structural properties of the data. The ProtT5 embedding space, with the smallest $\delta$ value $0.0418$ and ultrametricity value $0.1301$, reflects its tree-like structure with minimal deviation. The differences in hyperbolicity across these embeddings highlight how the geometric and topological characteristics of the PDB186 data influence the extent to which the space approximates a tree metric. In contrast, seqvec, with a $\delta$ value of $3.2018$, exhibits higher hyperbolicity, and similarly, Ultrametricity with a value of $16.6730$ for seqvec denotes that the embedding deviates from the tree-like structure.

% % \textcolor{red}{Prakash: Maybe we can keep these values side by side, depending on what we are comparing and what these results depict.}}

% % \textcolor{green}{
% The Neighbor Joining (NJ) scores differ substantially between Poincaré and Euclidean distance matrices. Euclidean distances yield significantly higher NJ average values (e.g., 318.63 for ProtT5 on PDB186) compared to Poincaré distances (567.03), indicating that Euclidean embeddings deviate more from an additive tree structure. The Poincaré ball naturally preserves hierarchical relationships through its hyperbolic geometry, yielding lower NJ scores that more closely approximate tree-like organization. This explains why Poincaré-based metrics show stronger tree-likeness even when Euclidean distances reveal the same underlying embedding structure.
% }


% \begin{table*}[ht]
% \centering
% \resizebox{0.8\textwidth}{!}{
% \begin{tabular}{cp{1.6cm}ccc|ccc|ccc}
% \toprule
% \multirow{2}{*}{Distance Matrix} 
% & \multirow{2}{*}{Embedding}
% & \multicolumn{3}{c}{Hyperbolicity} 
% & \multicolumn{3}{c}{Ultrametricity} 
% & \multicolumn{3}{c}{Neighbor-Joining (NJ)} \\
% \cmidrule{3-5} \cmidrule{6-8} \cmidrule{9-11}
% & 
% & $\delta_{\max}$ & $\delta_{\text{avg}}$ & $\delta_{\text{std}}$
% & $U_{\max}$ & $U_{\text{avg}}$ & $U_{\text{std}}$
% & $NJ_{\max}$ & $NJ_{\text{avg}}$ & $NJ_{\text{std}}$ \\
% \midrule

% \multirow{4}{*}{Poincaré}
% & SeqVec  & 0.2425 & 0.0110 & 0.0231 & 0.0937 & 0.0242 & 0.0208 & 74.1682 & 30.5558 & 6.2549 \\
% & ESM     & 0.1881 & 0.0183 & 0.0413 & \textbf{0.0405} & 0.0124 & 0.0068 & 62.8122 & 19.3127 & 11.7297 \\
% & TAPE    & \textbf{0.0794} & \textbf{0.0030} & 0.0063 & 0.0492 & \textbf{0.0083} & 0.0064 & \textbf{36.2088} & \textbf{16.3208} & 1.4938 \\
% & ProtT5  & 1.3882 & 0.0446 & 0.0957 & 2.9941 & 0.3038 & 0.2598 & 1023.8310 & 567.0299 & 24.3637 \\

% \midrule
% \multirow{4}{*}{Euclidean}
% & SeqVec  & 39.9819 & 1.6227 & 3.1682 & 56.8067 & 3.6570 & 4.1475 & 16313.2832 & 7276.7976 & 860.5404 \\
% & ESM     & 8.7726 & 0.8727 & 1.9674 & 1.9903 & 0.7193 & 0.3519 & 2902.5279 & 920.0283 & 559.1826 \\
% & TAPE    & 9.6066 & 0.4178 & 0.8754 & 6.1348 & 1.0126 & 0.7855 & 4672.6023 & 2123.9591 & 207.2407 \\
% & ProtT5  & \textbf{1.0945} & \textbf{0.0418} & 0.0905 & \textbf{0.9371} & \textbf{0.1300} & 0.1011 & \textbf{552.7800} & \textbf{318.6305} & 22.7829 \\

% \bottomrule
% \end{tabular}
% }
% \caption{Hyperbolicity, Ultrametricity, and Neighbor-Joining (NJ) statistics for the PDB186 dataset, grouped by the distance matrix used to compute embedding distances.}
% \label{tab_pdb186_distance_matrix_comparison}
% \end{table*}




% \begin{table*}[ht]
% \centering
% \resizebox{0.99\textwidth}{!}{
% \begin{tabular}{cp{1.4cm}ccc|ccc|ccc|ccc|ccc|ccc}
% \toprule
% \multirow{2}{*}{Distance Matrix} 
% & \multirow{2}{*}{Embedding}
% & \multicolumn{3}{c}{Hyperbolicity} 
% & \multicolumn{3}{c}{Ultrametricity} 
% & \multicolumn{3}{c}{Neighbor-Joining (NJ)} 
% & \multicolumn{3}{c}{Ricci curvature}
% & \multicolumn{3}{c}{Persistence entropy}
% & \multicolumn{3}{c}{Fractal dimension}
% \\
% \cmidrule{3-5} \cmidrule{6-8} \cmidrule{9-11} \cmidrule{12-14} \cmidrule{15-17} \cmidrule{18-20}
% & 
% & $\delta_{\max}$ & $\delta_{\text{avg}}$ & $\delta_{\text{std}}$
% & $U_{\max}$ & $U_{\text{avg}}$ & $U_{\text{std}}$
% & $NJ_{\max}$ & $NJ_{\text{avg}}$ & $NJ_{\text{std}}$ 
% & $ricci_{\max}$ & $ricci_{\text{avg}}$ & $ricci_{\text{std}}$ 
% & $pers_{\max}$ & $pers_{\text{avg}}$ & $pers_{\text{std}}$ 
% & $fractal_{\max}$ & $fractal_{\text{avg}}$ & $fractal_{\text{std}}$ 
% \\
% \midrule

% \multirow{4}{*}{Poincaré}
% & SeqVec  & 0.2425 & 0.0110 & 0.0231 & 0.0937 & 0.0242 & 0.0208 & 74.1682 & 30.5558 & 6.2549 & -0.0025 & -0.1712 & 0.1363 & 137.0438 & 19.0043 & 23.0151 & 886.2656 & 29.9867 & 89.2350 \\
% & ESM     & 0.1881 & 0.0183 & 0.0413 & \textbf{0.0405} & 0.0124 & 0.0068 & 62.8122 & 19.3127 & 11.7297 & -0.0000 & -0.0650 & 0.1893 & 7.4435 &  0.3970 &  1.4625 & inf &  inf &  nan              \\
% & TAPE    & \textbf{0.0794} & \textbf{0.0030} & 0.0063 & 0.0492 & \textbf{0.0083} & 0.0064 & \textbf{36.2088} & \textbf{16.3208} & 1.4938 & -0.0122 & -0.0642 & 0.0437 & 13.7103 & 7.3302 &  1.7907 & 58.8514 & 16.3534 & 8.9410     \\
% & ProtT5  & 1.3882 & 0.0446 & 0.0957 & 2.9941 & 0.3038 & 0.2598 & 1023.8310 & 567.0299 & 24.3637 & -0.0113 & -0.0652 & 0.0433 & 2.7212 &  1.1681 &  0.2538 & 64.7234 & 19.4370 & 12.336     \\

% \midrule
% \multirow{4}{*}{Euclidean}
% & SeqVec  & 39.9819 & 1.6227 & 3.1682 & 56.8067 & 3.6570 & 4.1475 & 16313.2832 & 7276.7976 & 860.5404 & -0.0025 & -0.1712 & 0.1363 & 137.0438 & 19.0043 & 23.0151 & 886.2656 & 29.9867 & 89.2350	\\
% & ESM     & 8.7726 & 0.8727 & 1.9674 & 1.9903 & 0.7193 & 0.3519 & 2902.5279 & 920.0283 & 559.1826 & -0.0000 & -0.0650 & 0.1893 & 7.4435 &  0.3970 &  1.4625 & inf & inf &  nan               \\
% & TAPE    & 9.6066 & 0.4178 & 0.8754 & 6.1348 & 1.0126 & 0.7855 & 4672.6023 & 2123.9591 & 207.2407 & -0.0122 & -0.0642 & 0.0437 & 13.7103 & 7.3302 &  1.7907 & 58.8514 & 16.3534 & 8.9410        \\
% & ProtT5  & \textbf{1.0945} & \textbf{0.0418} & 0.0905 & \textbf{0.9371} & \textbf{0.1300} & 0.1011 & \textbf{552.7800} & \textbf{318.6305} & 22.7829 & -0.0113 & -0.0652 & 0.0433 & 2.7212 &  1.1681 &  0.2538 & 64.7234 & 19.4370 & 12.3369       \\

% \bottomrule
% \end{tabular}
% }
% \caption{Hyperbolicity, Ultrametricity, Neighbor-Joining (NJ), Ricci, Persistence, and Fractal Statistics for the PDB186 dataset, grouped by the distance matrix used to compute embedding distances.}
% \label{tab_pdb186_distance_matrix_comparison_geometric}
% \end{table*}

\begin{table}[ht]
\centering
\LARGE
\resizebox{0.9\textwidth}{!}{
\begin{tabular}{lccccccc}
\toprule
\multirow{1}{*}{Dist Matrix} & Embedding
& $\boldsymbol{\delta_{\text{avg}}}$ 
& $\boldsymbol{U_{\text{avg}}}$ 
& $\boldsymbol{\mathrm{NJ}_{\text{avg}}}$ 
& $\boldsymbol{\mathrm{Ricci}_{\text{avg}}}$ 
& $\boldsymbol{\mathrm{PersEnt}_{\text{avg}}}$ 
& $\boldsymbol{\mathrm{FracDim}_{\text{avg}}}$ \\
\midrule
\multirow{4}{*}{Poincare}
& SeqVec  & 0.0110 & 0.0242 & 30.5558 & \textbf{-0.1712} & \textbf{19.0043} & 29.9867	\\
& ESM     & 0.0183 & 0.0124 & 19.3127 & -0.0650 & 0.3970 & inf        \\
& TAPE    & \textbf{0.0030} & \textbf{0.0083} & \textbf{16.3208} & -0.0642 & 7.3302 & \textbf{16.3534}    \\
& ProtT5  & 0.0446 & 0.3038 & 567.0299 & -0.0652 & 1.1681 & 19.4370     \\
\midrule
\multirow{4}{*}{Euclidean}
& SeqVec  & 1.6227 & 3.6570 & 7276.7976 & -0.1712 & 19.0043  & 29.9867	\\
& ESM     & 0.8727 & 0.7193 & 920.0283 & -0.0650 & 0.3970 & inf           \\
& TAPE    & 0.4178 & 1.0126 & 2123.9591 & -0.0642 & 7.3302 & 16.3534      \\
& ProtT5  & \textbf{0.0418} & \textbf{0.1300} & \textbf{318.6305} & -0.0652 & 1.1681 & 19.4370       \\
\bottomrule
\end{tabular}
}
\caption{Comparison of geometric and topological metrics across different embeddings for the PDB186 dataset, grouped by the distance matrix. The log ratio in the fractal dimension formula approaches zero, and its reciprocal explodes to infinity as a result of ESM embeddings placing numerous nearest neighbors at approximately identical distances, which results in the inf value.}
\label{tab_PDB186_embedding_metrics}
\end{table}


% \begin{table}[ht]
% \centering
% \resizebox{0.9\textwidth}{!}{
% \begin{tabular}{cp{0.6cm}cccc|cccc|cccc}
% \toprule
% \multirow{2}{*}{Embedding}

% & \multicolumn{5}{c}{$\boldsymbol{\mathrm{Ricci}}$ } 
% & \multicolumn{4}{c}{$\boldsymbol{\mathrm{PersEnt}}$ } 
% & \multicolumn{4}{c}{$\boldsymbol{\mathrm{FracDim}}$} \\
% \cmidrule{2-6} \cmidrule{7-10} \cmidrule{11-14}
% & $R_{\text{mean}}$ & $R_{\text{std}}$ & $R_{\text{min}}$ & $R_{\text{max}}$ & $R_{\text{neg ratio}}$
% & $PE_{\text{mean}}$ & $PE_{\text{std}}$ & $PE_{\text{max}}$ & $PE_{\text{entropy}}$
% & $FD_{\text{mean}}$ & $FD_{\text{std}}$ & $FD_{\text{max}}$ & $FD_{\text{median}}$ \\
% \midrule

% SeqVec  & & & & & & & & & & & \\
% ESM     & & & & & & & & & & & \\
% TAPE    & & & & & & & & & & & \\
% ProtT5  & & & & & & & & & & & \\

% \bottomrule
% \end{tabular}
% }
% \caption{Ricci, Persistence, and Fractal Statistics for the different embeddings for the PDB186 dataset.}
% \label{tab_pdb186_geometric_understanding}
% \end{table}


% \begin{table*}[ht]
% \centering
% \resizebox{\textwidth}{!}{
% \begin{tabular}{lccccc|cccc|cccc}
% \toprule
% \multirow{2}{*}{\textbf{Embedding}} 
% & \multicolumn{5}{c|}{\textbf{Ricci Curvature}} 
% & \multicolumn{4}{c|}{\textbf{Persistence (0D)}} 
% & \multicolumn{4}{c}{\textbf{Fractal Dimension}} \\

% \cmidrule(lr){2-6}
% \cmidrule(lr){7-10}
% \cmidrule(lr){11-14}

% & Mean $\downarrow$ & Std & Min & Max & Neg. Ratio $\uparrow$
% & Mean $\downarrow$ & Std & Max $\downarrow$ & Entropy
% & Mean $\downarrow$ & Std & Median $\downarrow$ & Max \\

% \midrule
% SeqVec 
% & \textbf{-0.1712} & 0.1363 & -0.5110 & -0.0025 & \textbf{1.000}
% & 19.0043 & 23.0151 & 137.0438 & 6.7301
% & 29.9867 & 89.2350 & \textbf{4.2842} & 886.2656 \\

% ESM 
% & -0.0650 & 0.1893 & -0.8174 & -0.0000 & 0.129
% & \textbf{0.3970} & 1.4625 & 7.4435 & \textbf{3.6828}
% & $\infty$ & NaN & $\infty$ & $\infty$ \\

% TAPE 
% & -0.0642 & \textbf{0.0437} & -0.2724 & -0.0122 & \textbf{1.000}
% & 7.3302 & 1.7907 & 13.7103 & 7.4904
% & \textbf{16.3534} & \textbf{8.9410} & 14.8255 & \textbf{58.8514} \\

% ProtT5 
% & -0.0652 & 0.0433 & -0.2273 & -0.0113 & \textbf{1.000}
% & 1.1681 & \textbf{0.2538} & \textbf{2.7212} & 7.4992
% & 19.4370 & 12.3369 & 16.1403 & 64.7234 \\

% \midrule
% \textbf{Tree-likeness ↑}
% & More neg. & -- & -- & -- & Higher
% & Lower & -- & Lower & --
% & Lower & -- & Lower & -- \\

% \bottomrule
% \end{tabular}
% }
% \caption{
% % Geometric and Topological Analysis of PDB186 Embeddings. Arrows indicate whether higher ($\uparrow$) or lower ($\downarrow$) values correspond to a stronger tree-like structure.
% % \textcolor{green}{
% Comprehensive geometric and topological analysis of PDB186 embeddings. Arrows indicate whether higher ($\uparrow$) or lower ($\downarrow$) values correspond to a stronger tree-like structure.}
% % }
% \label{tab_geometry_topology_PDB186_embeddings}
% \end{table*}


% #######Euclidean ##########################
% =================
% Formatted Table (All Metrics)
% =================
% Embedding  δ_avg    U_avg    NJ_avg   Ricci_avg PersEnt_avg FracDim_avg
% ----------------------------------------------------------------------------------------------------------------------------------------------------------------
% SeqVec     1.6227   3.6570   7276.7976 -0.1712  19.0043  29.9867
% ESM        0.8727   0.7193   920.0283 -0.0650  0.3970   inf
% TAPE       0.4178   1.0126   2123.9591 -0.0642  7.3302   16.3534
% ProtT5     0.0418   0.1300   318.6305 -0.0652  1.1681   19.4370
% ===================
% SeqVec
%   Ricci curvature: mean=-0.1712, std=0.1363, min=-0.5110, max=-0.0025, negative_ratio=1.0000
%   Persistence: mean=19.0043, std=23.0151, max=137.0438, entropy=6.7301
%   Fractal dimension: max=886.2656, mean=29.9867, std=89.2350, median=4.2842
% ESM
%   Ricci curvature: mean=-0.0650, std=0.1893, min=-0.8174, max=-0.0000, negative_ratio=0.1290
%   Persistence: mean=0.3970, std=1.4625, max=7.4435, entropy=3.6828
%   Fractal dimension: max=inf, mean=inf, std=nan, median=inf
% TAPE
%   Ricci curvature: mean=-0.0642, std=0.0437, min=-0.2724, max=-0.0122, negative_ratio=1.0000
%   Persistence: mean=7.3302, std=1.7907, max=13.7103, entropy=7.4904
%   Fractal dimension: max=58.8514, mean=16.3534, std=8.9410, median=14.8255
% ProtT5
%   Ricci curvature: mean=-0.0652, std=0.0433, min=-0.2273, max=-0.0113, negative_ratio=1.0000
%   Persistence: mean=1.1681, std=0.2538, max=2.7212, entropy=7.4992
%   Fractal dimension: max=64.7234, mean=19.4370, std=12.3369, median=16.1403

% ##########Euciedean###################################


% #####################Poincare###############
% Formatted Table (All Metrics)
% ================================================================================================================================================================
% Embedding  δ_avg    U_avg    NJ_avg   Ricci_avg PersEnt_avg FracDim_avg
% ----------------------------------------------------------------------------------------------------
% SeqVec     0.0110   0.0242   30.5558  -0.1712  19.0043  29.9867
% ESM        0.0183   0.0124   19.3127  -0.0650  0.3970   inf
% TAPE       0.0030   0.0083   16.3208  -0.0642  7.3302   16.3534
% ProtT5     0.0446   0.3038   567.0299 -0.0652  1.1681   19.4370
% ====================================================================================================
% SeqVec
%   Ricci curvature: mean=-0.1712, std=0.1363, min=-0.5110, max=-0.0025, negative_ratio=1.0000
%   Persistence: mean=19.0043, std=23.0151, max=137.0438, entropy=6.7301
%   Fractal dimension: max=886.2656, mean=29.9867, std=89.2350, median=4.2842
% ESM
%   Ricci curvature: mean=-0.0650, std=0.1893, min=-0.8174, max=-0.0000, negative_ratio=0.1290
%   Persistence: mean=0.3970, std=1.4625, max=7.4435, entropy=3.6828
%   Fractal dimension: max=inf, mean=inf, std=nan, median=inf
% TAPE
%   Ricci curvature: mean=-0.0642, std=0.0437, min=-0.2724, max=-0.0122, negative_ratio=1.0000
%   Persistence: mean=7.3302, std=1.7907, max=13.7103, entropy=7.4904
%   Fractal dimension: max=58.8514, mean=16.3534, std=8.9410, median=14.8255
% ProtT5
%   Ricci curvature: mean=-0.0652, std=0.0433, min=-0.2273, max=-0.0113, negative_ratio=1.0000
%   Persistence: mean=1.1681, std=0.2538, max=2.7212, entropy=7.4992
%   Fractal dimension: max=64.7234, mean=19.4370, std=12.3369, median=16.1403

% ###########################################




% \begin{table}[ht]
% \centering
% \resizebox{0.45\textwidth}{!}{
% \begin{tabular}{p{1.5cm}ccc|ccP{1cm}}
% \toprule
% \multirow{3}{1.5cm}{Embedding} 
% & \multicolumn{3}{c}{Hyperbolicity} 
% & \multicolumn{3}{c}{Ultrametricity} 
% % & \multirow{3}{3cm}{\centering Cophenetic Correlation Coefficient (CCC)} 
% \\
% \cmidrule{2-4} \cmidrule{5-7}
% & $\delta_{\text{max}}$ & $\delta_{\text{average}}$ & $\delta_{\text{std}}$ 
% & Max & Average & Std 
% % &
% \\
% \midrule
% Seqvec   & 116.7348 & 3.2018 & 6.9712 & 142.5040 & 16.6730 & 16.7692 
% % & \textbf{0.9812} 
% \\
% ESM2     & 6.0017   & 0.1913 & 0.4090 & 3.5801   & 0.4888  & 0.4281  
% % & 0.7604 
% \\
% TAPE     & 9.6066   & 0.4178 & 0.8754 & 6.1348   & 1.0116  & 0.7845  
% % & 0.8569 
% \\
% ProtT5   & \textbf{1.0945}   & \textbf{0.0418} & 0.0905 & \textbf{0.9371}   & \textbf{0.1301}  & 0.1011  
% % & 0.8742 
% \\
% \bottomrule
% \end{tabular}
% }
% \caption{Hyperbolicity, and Ultrametricity values (max, average, and standard deviation) for PDB186 dataset.}
% \label{tab_pdb186_hyperbolicity_ultramtricity}
% \end{table}


% \begin{table}[h!]
%     \centering
%     \begin{tabular}{cccc}
%     \toprule
%     Embedding & $\delta_{max}$ & $\delta_{Average}$ & $\delta_{Std.}$ \\
%     \midrule \midrule
%        Seqvec & 25.1741 & 0.5925 & 1.1635 \\
%        ESM2 & 21.0830 & 1.0928 & 2.1575 \\
%        % TAPE & 11.7150 & 0.4973 & 1.1015 \\
%        TAPE & 11.7450 & 0.5905 & 1.1867 \\
%        ProtT5 & 2.9494 & 0.1013 & 0.2171 \\
%        \bottomrule
%     \end{tabular}
%     \caption{Hyperbolicity $\delta$ values (average and standard deviation) for\textbf{Lung Cancer} datasets.}
%     \label{tbl_delta}
% \end{table}

% \begin{table}[h!]
%     \centering
%     \begin{tabular}{cccc}
%     \toprule
%     Embedding & $max$ & $Average$ & $Std.$\\
%     \midrule \midrule
%        Seqvec & 13.5444 &  1.3015 & 1.1148 \\
%        ESM2 & 8.5433 &  1.5102  & 1.1435 \\
%        TAPE & 7.0010 &  1.2213  & 0.9452 \\
%        ProtT5 & 1.9276 &  0.2408  & 0.1984 \\
       
%        \bottomrule
%     \end{tabular}
%     \caption{Ultrametricity values (max, and average) for\textbf{Lung Cancer} datasets.}
%     \label{tbl_delta}
% \end{table}

\subsubsection{Lung Cancer dataset}
% The hyperbolicity differences in Table~\ref{tab_lung_cancer_distance_matrix_comparison} reflect the structural properties of the Lung Cancer embeddings. ProtT5 shows the most tree-like geometry, with the lowest $\delta$ (0.1013) and ultrametricity (0.2408), indicating minimal deviation from a hierarchical structure. In contrast, SeqVec—with a higher $\delta$ (0.5925) and ultrametricity (1.2458)—deviates more strongly from a tree metric. These variations highlight how each embedding’s geometry influences its alignment with hierarchical structure.

The Lung Cancer results show that geometric behavior varies strongly across embeddings and distance metrics. As shown in Table~\ref{tab_lung_cancer_embedding_metrics}, in the Poincaré space, ESM2 is the most hierarchical, with the lowest hyperbolicity (0.0016), ultrametricity (0.0022), and NJ deviation (40.29), as well as the most negative Ricci curvature. TAPE shows the lowest fractal dimension (18.57), indicating a more compact manifold, while SeqVec exhibits the highest persistence entropy, reflecting greater multiscale variability. ProtT5 shows weak hierarchy in Poincaré space but becomes the most tree-like under Euclidean distances, with the lowest $\delta$ (0.1010), $U$ (0.2413), and $NJ$ deviation (2296.54). Overall, ESM2 and TAPE demonstrate stronger hyperbolic structure, whereas ProtT5 appears more separable in Euclidean geometry.
% The variation in hyperbolicity values for the Lung Cancer dataset, as shown in Table~\ref{tab_lung_cancer_distance_matrix_comparison}, can be attributed to the underlying structural properties of the data. The ProtT5 embedding space, with the smallest $\delta$ value of $0.1013$ and an ultrametricity value of $0.2408$, reflects its tree-like structure with minimal deviation. The differences in hyperbolicity across these embeddings highlight how the geometric and topological characteristics of the data influence the extent to which the space approximates a tree metric. In contrast, the seqvec, with a $\delta$ value of $0.5925$, exhibits higher hyperbolicity. Similarly, Ultrametricity, with a value of $1.2458$ for seqvec, denotes that the embedding deviates from a tree-like structure.


% \begin{table*}[ht]
% \centering
% \resizebox{0.8\textwidth}{!}{
% \begin{tabular}{cp{1.6cm}ccc|ccc|ccc}
% \toprule
% \multirow{2}{*}{Distance Matrix} 
% & \multirow{2}{*}{Embedding}
% & \multicolumn{3}{c}{Hyperbolicity} 
% & \multicolumn{3}{c}{Ultrametricity} 
% & \multicolumn{3}{c}{Neighbor-Joining (NJ)} \\
% \cmidrule{3-5} \cmidrule{6-8} \cmidrule{9-11}
% & 
% & $\delta_{\max}$ & $\delta_{\text{avg}}$ & $\delta_{\text{std}}$
% & $U_{\max}$ & $U_{\text{avg}}$ & $U_{\text{std}}$
% & $NJ_{\max}$ & $NJ_{\text{avg}}$ & $NJ_{\text{std}}$ \\
% \midrule

% \multirow{4}{*}{Poincaré}
% & SeqVec & 0.0793 & 0.0026 & 0.0053 & 0.0369 & 0.0069 & 0.0056 & 224.0834 & 118.3239 & 6.8652		\\
% & ESM &  \textbf{0.0329} & \textbf{0.0016} & 0.0032 & \textbf{0.0162} & \textbf{0.0022} & 0.0018 & \textbf{74.9932} &  \textbf{40.2858} &  4.0370       \\
% & TAPE & 0.0487 & 0.0040 & 0.0080 & 0.0373 & 0.0084 & 0.0064 & 138.4172 & 72.4654 &  9.7122       \\
% & ProtT5 & 0.6608 & 0.0305 & 0.0614 & 0.3996 & 0.0759 & 0.0604 & 1294.4451 & 715.5995 & 73.8501   \\

% \midrule
% \multirow{4}{*}{Euclidean}
% & SeqVec & 16.6321 & 0.5904 & 1.1575 & 8.8284 & 1.2987 & 1.0915 & 51448.2682 & 28674.5853 & 1470.6535 \\
% & ESM    & 22.6119 & 1.0897 & 2.1588 & 7.5962 & 1.4986 & 1.1505 & 52579.6242 & 27513.1399 & 2739.3525 \\
% & TAPE   & 7.0102  & 0.5944 & 1.1770 & 6.2529 & 1.2251 & 0.9505 & 23307.1818 & 10738.2480 & 1437.5647 \\
% & ProtT5 & \textbf{2.6989}  & \textbf{0.1010} & 0.2241 & \textbf{1.7052} & \textbf{0.2413} & 0.2022 & \textbf{5254.1620}  & \textbf{2296.5416}  & 263.0430  \\

% \bottomrule
% \end{tabular}
% }
% \caption{Hyperbolicity, Ultrametricity, and Neighbor-Joining (NJ) statistics for the Lung Cancer dataset, grouped by the distance matrix used to compute embedding distances.}
% \label{tab_lung_cancer_distance_matrix_comparison}
% \end{table*}

% \begin{table*}[ht]
% \centering
% \resizebox{0.99\textwidth}{!}{
% \begin{tabular}{cp{1.4cm}ccc|ccc|ccc|ccc|ccc|ccc}
% \toprule
% \multirow{2}{*}{Distance Matrix} 
% & \multirow{2}{*}{Embedding}
% & \multicolumn{3}{c}{Hyperbolicity} 
% & \multicolumn{3}{c}{Ultrametricity} 
% & \multicolumn{3}{c}{Neighbor-Joining (NJ)} 
% & \multicolumn{3}{c}{Ricci curvature}
% & \multicolumn{3}{c}{Persistence entropy}
% & \multicolumn{3}{c}{Fractal dimension}
% \\
% \cmidrule{3-5} \cmidrule{6-8} \cmidrule{9-11} \cmidrule{12-14} \cmidrule{15-17} \cmidrule{18-20}
% & 
% & $\delta_{\max}$ & $\delta_{\text{avg}}$ & $\delta_{\text{std}}$
% & $U_{\max}$ & $U_{\text{avg}}$ & $U_{\text{std}}$
% & $NJ_{\max}$ & $NJ_{\text{avg}}$ & $NJ_{\text{std}}$ 
% & $ricci_{\max}$ & $ricci_{\text{avg}}$ & $ricci_{\text{std}}$ 
% & $pers_{\max}$ & $pers_{\text{avg}}$ & $pers_{\text{std}}$ 
% & $fractal_{\max}$ & $fractal_{\text{avg}}$ & $fractal_{\text{std}}$ 
% \\
% \midrule

% \multirow{4}{*}{Poincaré}
% & SeqVec & 0.0793 & 0.0026 & 0.0053 & 0.0369 & 0.0069 & 0.0056 & 224.0834 & 118.3239 & 6.8652 & -0.0055 & -0.0637 & 0.1004 & 45.5655 & 22.4044 & 7.1787 & 124.1813 & 32.0288 & 18.1777	\\
% & ESM &  \textbf{0.0329} & \textbf{0.0016} & 0.0032 & \textbf{0.0162} & \textbf{0.0022} & 0.0018 & \textbf{74.9932} &  \textbf{40.2858} &  4.0370       & -0.0005 & -0.0849 & 0.2089 & 27.6253 & 20.7693 & 6.2660 & 592.5030 & 32.9002 & 27.0752  \\
% & TAPE & 0.0487 & 0.0040 & 0.0080 & 0.0373 & 0.0084 & 0.0064 & 138.4172 & 72.4654 &  9.7122     & -0.0093 & -0.0653 & 0.0651 & 11.5238 & 6.4817 &  1.4347 & 76.3705 & 18.5660 & 9.6267    \\
% & ProtT5 & 0.6608 & 0.0305 & 0.0614 & 0.3996 & 0.0759 & 0.0604 & 1294.4451 & 715.5995 & 73.8501   & -0.0056 & -0.0621 & 0.0620 & 3.2069 &  1.4096 &  0.3807 & 111.3154 & 20.9063 & 11.7847  \\

% \midrule
% \multirow{4}{*}{Euclidean}
% & SeqVec & 16.6321 & 0.5904 & 1.1575 & 8.8284 & 1.2987 & 1.0915 & 51448.2682 & 28674.5853 & 1470.6535 & -0.0055 & -0.0637 & 0.1004 & 45.5655 & 22.4044 & 7.1787 & 124.1813 & 32.0288 & 18.1777	\\
% & ESM    & 22.6119 & 1.0897 & 2.1588 & 7.5962 & 1.4986 & 1.1505 & 52579.6242 & 27513.1399 & 2739.3525 & -0.0005 & -0.0849 & 0.2089 & 27.6253 & 20.7693 & 6.2660 & 592.5030 & 32.9002 & 27.0752    \\
% & TAPE   & 7.0102  & 0.5944 & 1.1770 & 6.2529 & 1.2251 & 0.9505 & 23307.1818 & 10738.2480 & 1437.5647 & -0.0093 & -0.0653 & 0.0651 & 11.5238 & 6.4817 &  1.4347 & 76.3705 & 18.5660 & 9.6267      \\
% & ProtT5 & \textbf{2.6989}  & \textbf{0.1010} & 0.2241 & \textbf{1.7052} & \textbf{0.2413} & 0.2022 & \textbf{5254.1620}  & \textbf{2296.5416}  & 263.0430  & -0.0056 & -0.0621 & 0.0620 & 3.2069 &  1.4096 &  0.3807 & 111.3154 & 20.9063 & 11.7847    \\

% \bottomrule
% \end{tabular}
% }
% \caption{Hyperbolicity, Ultrametricity, Neighbor-Joining (NJ), Ricci, Persistence, and Fractal statistics for the Lung Cancer dataset, grouped by the distance matrix used to compute embedding distances.}
% \label{tab_lung_cancer_geometric_comparison}
% \end{table*}

\begin{table}[ht]
\centering
\LARGE
\resizebox{0.9\textwidth}{!}{
\begin{tabular}{lccccccc}
\toprule
\multirow{1}{*}{Dist Matrix} & Embedding
& $\boldsymbol{\delta_{\text{avg}}}$ 
& $\boldsymbol{U_{\text{avg}}}$ 
& $\boldsymbol{\mathrm{NJ}_{\text{avg}}}$ 
& $\boldsymbol{\mathrm{Ricci}_{\text{avg}}}$ 
& $\boldsymbol{\mathrm{PersEnt}_{\text{avg}}}$ 
& $\boldsymbol{\mathrm{FracDim}_{\text{avg}}}$ \\
\midrule
\multirow{4}{*}{Poincare}
& SeqVec & 0.0026 & 0.0069 & 118.3239 	& -0.0637 & \textbf{22.4044} & 32.0288	\\
& ESM &  \textbf{0.0016} & \textbf{0.0022} &  \textbf{40.2858} & \textbf{-0.0849} & 20.7693 & 32.9002 \\
& TAPE & 0.0040  & 0.0084 &  72.4654   & -0.0653 & 6.4817 &  \textbf{18.5660} \\
& ProtT5 & 0.0305 & 0.0759 & 715.5995  & -0.0621 & 1.4096 &  20.9063 \\

\midrule
\multirow{4}{*}{Euclidean}
& SeqVec & 0.5904 & 1.2987 & 28674.5853 & -0.0637 & \textbf{22.4044} & 32.0288 \\
& ESM    & 1.0897 & 1.4986 & 27513.1399 &\textbf{ -0.0849} & 20.7693 & 32.9002 \\
& TAPE   & 0.5944 & 1.2251 & 10738.2480 & -0.0653 & 6.4817 &  \textbf{18.5660} \\
& ProtT5 & \textbf{0.1010} & \textbf{0.2413} & \textbf{2296.5416} & -0.0621 & 1.4096 &  20.9063 \\
\bottomrule
\end{tabular}
}
\caption{Comparison of geometric and topological metrics across different embeddings for the Lung Cancer dataset, grouped by the distance matrix.}
\label{tab_lung_cancer_embedding_metrics}
\end{table}

% \begin{table*}[ht]
% \centering
% \resizebox{\textwidth}{!}{
% \begin{tabular}{lccccc|cccc|cccc}
% \toprule
% \multirow{2}{*}{\textbf{Embedding}} 
% & \multicolumn{5}{c|}{\textbf{Ricci Curvature}} 
% & \multicolumn{4}{c|}{\textbf{Persistence (0D)}} 
% & \multicolumn{4}{c}{\textbf{Fractal Dimension}} \\

% \cmidrule(lr){2-6}
% \cmidrule(lr){7-10}
% \cmidrule(lr){11-14}

% & Mean $\downarrow$ & Std & Min & Max & Neg. Ratio $\uparrow$
% & Mean $\downarrow$ & Std & Max $\downarrow$ & Entropy
% & Mean $\downarrow$ & Std & Median $\downarrow$ & Max \\

% \midrule
% SeqVec 
% & -0.0637 & 0.1004 & -0.6821 & -0.0055 & \textbf{1.000}
% & 22.4044 & 7.1787 & 45.5655 & 9.7273
% & 32.0288 & 18.1777 & 32.0460 & 124.1813 \\

% ESM 
% & \textbf{-0.0849} & 0.2089 & \textbf{-1.5505} & -0.0005 & \textbf{1.000}
% & 20.7693 & 6.2660 & 27.6253 & 9.7134
% & 32.9002 & 27.0752 & 31.3154 & 592.5030 \\

% TAPE 
% & -0.0653 & \textbf{0.0651} & -0.5063 & -0.0093 & \textbf{1.000}
% & 6.4817 & 1.4347 & 11.5238 & \textbf{9.7730}
% & \textbf{18.5660} & \textbf{9.6267} & \textbf{17.9168} & \textbf{76.3705} \\

% ProtT5 
% & -0.0621 & 0.0620 & -0.4296 & -0.0056 & \textbf{1.000}
% & \textbf{1.4096} & \textbf{0.3807} & \textbf{3.2069} & 9.7620
% & 20.9063 & 11.7847 & 19.9629 & 111.3154 \\

% \midrule
% \textbf{Tree-likeness ↑}
% & More neg. & -- & -- & -- & Higher
% & Lower & -- & Lower & --
% & Lower & -- & Lower & -- \\

% \bottomrule
% \end{tabular}
% }
% \caption{Geometric and Topological Analysis of Lung Cancer Protein Embeddings.
% Arrows indicate whether higher ($\uparrow$) or lower ($\downarrow$) values correspond to stronger tree-like structure.}
% \label{tab:lung_geometry_topology}
% \end{table*}



% \begin{table}[ht]
% \centering
% \resizebox{0.45\textwidth}{!}{
% \begin{tabular}{p{1.5cm}ccc|ccP{1cm}}
% \toprule
% \multirow{3}{1.5cm}{Embedding} 
% & \multicolumn{3}{c}{Hyperbolicity} 
% & \multicolumn{3}{c}{Ultrametricity} 
% % & \multirow{3}{3cm}{\centering Cophenetic Correlation Coefficient (CCC)} 
% \\
% \cmidrule{2-4} \cmidrule{5-7}
% & $\delta_{\text{max}}$ & $\delta_{\text{average}}$ & $\delta_{\text{std}}$ 
% & Max & Average & Std 
% % & 
% \\
% \midrule
% Seqvec   & 25.1741 & 0.5925 & 1.1635 & 13.5444 & 1.2458 & 1.1148 
% % & \textbf{0.9573} 
% \\
% ESM2     & 21.0830 & 0.2982 & 2.1575 & 8.5433  & 1.5102 & 1.1435 
% % & 0.8526 
% \\
% TAPE     & 11.7450 & 0.5905 & 1.1867 & 7.0010  & 1.2213 & 0.9452 
% % & 0.6255 
% \\
% ProtT5   & \textbf{2.9494}  & \textbf{0.1013} & 0.2171 & \textbf{1.9276}  & \textbf{0.2408} & 0.1984 
% % &  0.8699 
% \\
% \bottomrule
% \end{tabular}
% }
% \caption{Hyperbolicity, and Ultrametricity values (max, average, and standard deviation) for Lung Cancer datasets.}
% \label{tab_lungcancer_metrics}
% \end{table}


% \begin{table}[h!]
%     \centering
%     \begin{tabular}{cccc}
%     \toprule
%     Embedding & $\delta_{max}$ & $\delta_{Average}$ & $\delta_{Std.}$ \\
%     \midrule \midrule
%        Seqvec & 25.9145 & 0.5882 & 1.1547 \\
%        ESM2 & 20.3740 & 1.0981 & 2.1559 \\
%        TAPE & 10.8329 &  0.6106 & 1.2126 \\
%        ProtT5 & 3.1166 & 0.1049 &  0.2275 \\
       
%        \bottomrule
%     \end{tabular}
%     \caption{Hyperbolicity $\delta$ values (average and standard deviation) for\textbf{Breast Cancer} datasets.}
%     \label{tbl_delta}
% \end{table}


% \begin{table}[h!]
%     \centering
%     \begin{tabular}{cccc}
%     \toprule
%     Embedding & $max$ & $Average$ & $Std.$\\
%     \midrule \midrule
%        Seqvec & 13.2806 &  1.3272 & 1.1422 \\
%        ESM2 & 9.1135 &  1.5185  & 1.1596 \\
%        TAPE & 7.3417 & 1.2469  & 0.9717 \\
%        ProtT5 & 1.8990 &  0.2466  & 0.2045 \\
       
%        \bottomrule
%     \end{tabular}
%     \caption{Ultrametricity values (max,  average, and fraction) for\textbf{Breast Cancer} datasets.}
%     \label{tbl_delta}
% \end{table}

\subsubsection{Breast Cancer dataset}
The Breast Cancer results in Table~\ref{tab_breast_cancer_embedding_metrics} show consistent geometric differences across embeddings. In the Poincaré space, ESM2 is the most hierarchical, achieving the lowest hyperbolicity (0.0017), ultrametricity (0.0023), and NJ deviation (42.74), along with the most negative Ricci curvature. Here also, TAPE shows the lowest fractal dimension (18.12), indicating a compact embedding manifold, while SeqVec exhibits the highest persistence entropy (22.23), reflecting more topological diversity. 
With significantly higher ultrametricity (0.0758) and NJ values (752.72), ProtT5 exhibits the weakest Poincaré hierarchy. Similar to patterns seen in previous datasets, ProtT5 becomes the most tree-like under Euclidean distances, exhibiting the lowest $\delta$ (0.1073), $U$ (0.2458), and $NJ$ deviation (2487.96). Overall, we can conclude here that ESM2 and TAPE exhibit stronger hyperbolic structure, whereas ProtT5 appears more separable in Euclidean space, and SeqVec remains more topologically fragmented.

% The variation in hyperbolicity values for the Breast Cancer dataset, as shown in Table~\ref{tab_breast_cancer_distance_matrix_comparison}, can be attributed to the underlying structural properties. The ProtT5 embedding space, with the smallest $\delta$ value with an average of $0.1049$ and an ultrametricity value of $0.2466$, reflects its tree-like structure with minimal deviation. 
% The differences in hyperbolicity across these embeddings highlight how the geometric and topological characteristics of the data influence the extent to which the space approximates a tree metric. 
% In contrast, the seqvec, with a $\delta$ value of $0.5882$, exhibits higher hyperbolicity. Similarly, Ultrametricity, with a value of $1.3272$ for seqvec, denotes that the embedding deviates from a tree-like structure.

% \begin{table*}[ht]
% \centering
% \resizebox{0.8\textwidth}{!}{
% \begin{tabular}{cp{1.6cm}ccc|ccc|ccc}
% \toprule
% \multirow{2}{*}{Distance Matrix} 
% & \multirow{2}{*}{Embedding}
% & \multicolumn{3}{c}{Hyperbolicity} 
% & \multicolumn{3}{c}{Ultrametricity} 
% & \multicolumn{3}{c}{Neighbor-Joining (NJ)} \\
% \cmidrule{3-5} \cmidrule{6-8} \cmidrule{9-11}
% & 
% & $\delta_{\max}$ & $\delta_{\text{avg}}$ & $\delta_{\text{std}}$
% & $U_{\max}$ & $U_{\text{avg}}$ & $U_{\text{std}}$
% & $NJ_{\max}$ & $NJ_{\text{avg}}$ & $NJ_{\text{std}}$ \\
% \midrule

% \multirow{4}{*}{Poincaré}
% & SeqVec &  0.0697 & 0.0025 & 0.0051  & 0.0371 & 0.0073 & 0.0059   &  228.8500  & 124.5329  & 7.1761 \\
% & ESM    &  \textbf{0.0296} & \textbf{0.0017} & 0.0033  &  \textbf{0.0155} & \textbf{0.0023} & 0.0019 &  \textbf{89.7679}   & \textbf{42.7396}  & 4.3141  \\
% & TAPE   &  0.0567 & 0.0040 & 0.0080  &  0.0362 & 0.0087 & 0.0065 &  146.2524  & 76.5285   & 10.0934 \\
% & ProtT5 &  0.4832 & 0.0306 & 0.0624  &  0.4166 & 0.0758 & 0.0606 &  1404.9604 & 752.7241  & 78.6217 \\


% \midrule
% \multirow{4}{*}{Euclidean}
% & SeqVec & 25.1794 & 0.5804 & 1.1569 & 9.7149 & 1.3425 & 1.1646 & 52250.6287 & 30408.3304 & 1559.0493 \\
% & ESM    & 16.4589 & 1.1336 & 2.1959 & 9.0446 & 1.5309 & 1.1629 & 54383.9908 & 28876.8367 & 2875.3675 \\
% & TAPE   & 8.5254  & 0.6024 & 1.2136 & 6.0660 & 1.2547 & 0.9680 & 23754.6981 & 11401.2842 & 1525.5165 \\ 
% & ProtT5 & \textbf{2.5025}  & \textbf{0.1073} & 0.2356 & \textbf{1.6354} & \textbf{0.2458} & 0.2044 & \textbf{6825.7746}  & \textbf{2487.9561}  & 294.8333 \\

% \bottomrule
% \end{tabular}
% }
% \caption{Hyperbolicity, Ultrametricity, and Neighbor-Joining (NJ) statistics for the Breast Cancer dataset, grouped by the distance matrix used to compute embedding distances.}
% \label{tab_breast_cancer_distance_matrix_comparison}
% \end{table*}


% \begin{table*}[ht]
% \centering
% \resizebox{0.99\textwidth}{!}{
% \begin{tabular}{cp{1.4cm}ccc|ccc|ccc|ccc|ccc|ccc}
% \toprule
% \multirow{2}{*}{Distance Matrix} 
% & \multirow{2}{*}{Embedding}
% & \multicolumn{3}{c}{Hyperbolicity} 
% & \multicolumn{3}{c}{Ultrametricity} 
% & \multicolumn{3}{c}{Neighbor-Joining (NJ)} 
% & \multicolumn{3}{c}{Ricci curvature}
% & \multicolumn{3}{c}{Persistence entropy}
% & \multicolumn{3}{c}{Fractal dimension}
% \\
% \cmidrule{3-5} \cmidrule{6-8} \cmidrule{9-11} \cmidrule{12-14} \cmidrule{15-17} \cmidrule{18-20}
% & 
% & $\delta_{\max}$ & $\delta_{\text{avg}}$ & $\delta_{\text{std}}$
% & $U_{\max}$ & $U_{\text{avg}}$ & $U_{\text{std}}$
% & $NJ_{\max}$ & $NJ_{\text{avg}}$ & $NJ_{\text{std}}$ 
% & $ricci_{\max}$ & $ricci_{\text{avg}}$ & $ricci_{\text{std}}$ 
% & $pers_{\max}$ & $pers_{\text{avg}}$ & $pers_{\text{std}}$ 
% & $fractal_{\max}$ & $fractal_{\text{avg}}$ & $fractal_{\text{std}}$ 
% \\
% \midrule

% \multirow{4}{*}{Poincaré}
% & SeqVec &  0.0697 & 0.0025 & 0.0051  & 0.0371 & 0.0073 & 0.0059   &  228.8500  & 124.5329  & 7.1761   & -0.0072 & -0.0675 & 0.1014 & 46.9667 & 22.2295 & 7.3156 & 121.0712 & 30.0684 & 17.5609	\\
% & ESM    &  \textbf{0.0296} & \textbf{0.0017} & 0.0033  &  \textbf{0.0155} & \textbf{0.0023} & 0.0019 &  \textbf{89.7679}   & \textbf{42.7396}  & 4.3141    & -0.0064 & -0.0748 & 0.1415 & 27.6253 & 20.4887 & 6.1579 & 136.4443 & 31.0122 & 19.3224 \\
% & TAPE   &  0.0567 & 0.0040 & 0.0080  &  0.0362 & 0.0087 & 0.0065 &  146.2524  & 76.5285   & 10.0934   & -0.0093 & -0.0690 & 0.0693 & 17.0604 & 6.4066 &  1.5882 & 76.3705 & 18.1227 & 9.6293   \\
% & ProtT5 &  0.4832 & 0.0306 & 0.0624  &  0.4166 & 0.0758 & 0.0606 &  1404.9604 & 752.7241  & 78.6217   & -0.0056 & -0.0653 & 0.0653 & 3.6860 &  1.4063 &  0.4012 & 111.3154 & 20.4609 & 11.8854 \\


% \midrule
% \multirow{4}{*}{Euclidean}
% & SeqVec & 25.1794 & 0.5804 & 1.1569 & 9.7149 & 1.3425 & 1.1646 & 52250.6287 & 30408.3304 & 1559.0493  & -0.0072 & -0.0675 & 0.1014 & 46.9667 & 22.2295 & 7.3156 & 121.0712 & 30.0684 & 17.5609	\\
% & ESM    & 16.4589 & 1.1336 & 2.1959 & 9.0446 & 1.5309 & 1.1629 & 54383.9908 & 28876.8367 & 2875.3675  & -0.0064 & -0.0748 & 0.1415 & 27.6253 & 20.4887 & 6.1579 & 136.4443 & 31.0122 & 19.3224  \\
% & TAPE   & 8.5254  & 0.6024 & 1.2136 & 6.0660 & 1.2547 & 0.9680 & 23754.6981 & 11401.2842 & 1525.5165  & -0.0093 & -0.0690 & 0.0693 & 17.0604 & 6.4066 &  1.5882 & 76.3705 & 18.1227 & 9.6293    \\
% & ProtT5 & \textbf{2.5025}  & \textbf{0.1073} & 0.2356 & \textbf{1.6354} & \textbf{0.2458} & 0.2044 & \textbf{6825.7746}  & \textbf{2487.9561}  & 294.8333 & -0.0056 & -0.0653 & 0.0653 & 3.6860 &  1.4063 &  0.4012 & 111.3154 & 20.4609 & 11.8854  \\

% \bottomrule
% \end{tabular}
% }
% \caption{Hyperbolicity, Ultrametricity, Neighbor-Joining (NJ), Ricci, Persistence, and Fractal statistics for the Breast Cancer dataset, grouped by the distance matrix used to compute embedding distances.}
% \label{tab_breast_cancer_geometric_comparison}
% \end{table*}


\begin{table}[ht]
\centering
\LARGE
\resizebox{0.9\textwidth}{!}{
\begin{tabular}{lccccccc}
\toprule
\multirow{1}{*}{Dist Matrix} & Embedding
& $\boldsymbol{\delta_{\text{avg}}}$ 
& $\boldsymbol{U_{\text{avg}}}$ 
& $\boldsymbol{\mathrm{NJ}_{\text{avg}}}$ 
& $\boldsymbol{\mathrm{Ricci}_{\text{avg}}}$ 
& $\boldsymbol{\mathrm{PersEnt}_{\text{avg}}}$ 
& $\boldsymbol{\mathrm{FracDim}_{\text{avg}}}$ \\
\midrule
\multirow{4}{*}{Poincare}
& SeqVec &  0.0025 & 0.0073 & 124.5329  & -0.0675 & 22.2295 & 30.0684 \\
& ESM    &  \textbf{0.0017} & \textbf{0.0023} & \textbf{42.7396} & -0.0748 & 20.4887 & 31.0122 \\
& TAPE   &  0.0040 & 0.0087 & 76.5285  & -0.0690 & 6.4066 &  18.1227 \\
& ProtT5 &  0.0306 & 0.0758 & 752.7241 & -0.0653 & 1.4063 &  20.4609 \\

\midrule
\multirow{4}{*}{Euclidean}
& SeqVec & 0.5804 & 1.3425 & 30408.3304 & -0.0675 & \textbf{22.2295} & 30.0684 \\
& ESM    & 1.1336 & 1.5309 & 28876.8367 & \textbf{-0.0748} & 20.4887 & 31.0122 \\
& TAPE   & 0.6024 & 1.2547 & 11401.2842 & -0.0690 & 6.4066 &  \textbf{18.1227} \\
& ProtT5 & \textbf{0.1073} & \textbf{0.2458} & \textbf{2487.9561} & -0.0653 & 1.4063 &  20.4609 \\
\bottomrule
\end{tabular}
}
\caption{Comparison of geometric and topological metrics across different embeddings for the Breast Cancer dataset, grouped by the distance matrix.}
\label{tab_breast_cancer_embedding_metrics}
\end{table}

\begin{remark}
    Because Ricci curvature, persistent homology, and fractal dimension are calculated directly on the embeddings rather than the complete pairwise distance matrix, they remain constant across the Euclidean and Poincaré distance matrices.
\end{remark}

% \begin{table*}[ht]
% \centering
% \resizebox{\textwidth}{!}{
% \begin{tabular}{lccccc|cccc|cccc}
% \toprule
% \multirow{2}{*}{\textbf{Embedding}} 
% & \multicolumn{5}{c|}{\textbf{Ricci Curvature}} 
% & \multicolumn{4}{c|}{\textbf{Persistence (0D)}} 
% & \multicolumn{4}{c}{\textbf{Fractal Dimension}} \\

% \cmidrule(lr){2-6}
% \cmidrule(lr){7-10}
% \cmidrule(lr){11-14}

% & Mean $\downarrow$ & Std & Min & Max & Neg. Ratio $\uparrow$
% & Mean $\downarrow$ & Std & Max $\downarrow$ & Entropy
% & Mean $\downarrow$ & Std & Median $\downarrow$ & Max \\

% \midrule
% SeqVec 
% & -0.0675 & 0.1014 & -0.6443 & -0.0072 & \textbf{1.000}
% & 22.2295 & 7.3156 & 46.9667 & 9.7976
% & 30.0684 & 17.5609 & 30.2730 & 121.0712 \\

% ESM 
% & \textbf{-0.0748} & 0.1415 & \textbf{-0.8240} & -0.0064 & \textbf{1.000}
% & 20.4887 & 6.1579 & 27.6253 & 9.7929
% & 31.0122 & 19.3224 & 30.1939 & 136.4443 \\

% TAPE 
% & -0.0690 & \textbf{0.0693} & -0.5118 & \textbf{-0.0093} & \textbf{1.000}
% & 6.4066 & 1.5882 & 17.0604 & \textbf{9.8381}
% & \textbf{18.1227} & \textbf{9.6293} & \textbf{17.6672} & \textbf{76.3705} \\

% ProtT5 
% & -0.0653 & 0.0653 & -0.4414 & -0.0056 & \textbf{1.000}
% & \textbf{1.4063} & \textbf{0.4012} & \textbf{3.6860} & 9.8311
% & 20.4609 & 11.8854 & 19.4204 & 111.3154 \\

% \midrule
% \textbf{Tree-likeness ↑}
% & More neg. & -- & -- & -- & Higher
% & Lower & -- & Lower & --
% & Lower & -- & Lower & -- \\

% \bottomrule
% \end{tabular}
% }
% \caption{Geometric and Topological Analysis of Breast Cancer Protein Embeddings.
% Arrows indicate whether higher ($\uparrow$) or lower ($\downarrow$) values correspond to a stronger tree-like structure.}
% \label{tab_breast_geometry_topology}
% \end{table*}


% \begin{table}[ht]
% \centering
% \resizebox{0.45\textwidth}{!}{
% \begin{tabular}{p{1.5cm}ccc|ccP{1cm}}
% \toprule
% \multirow{3}{1.5cm}{Embedding} 
% & \multicolumn{3}{c}{Hyperbolicity} 
% & \multicolumn{3}{c}{Ultrametricity} 
% % & \multirow{3}{3cm}{\centering Cophenetic Correlation Coefficient (CCC)} 
% \\
% \cmidrule{2-4} \cmidrule{5-7}
% & $\delta_{\text{max}}$ & $\delta_{\text{average}}$ & $\delta_{\text{std}}$ 
% & Max & Average & Std 
% % & 
% \\
% \midrule
% Seqvec   & 25.9145 & 0.5882 & 1.1547 & 13.2806 & 1.3272 & 1.1422 
% % & \textbf{0.9600} 
% \\
% ESM2     & 20.3740 & 1.0981 & 2.1559 & 9.1135  & 1.5185 & 1.1596 
% % & 0.8477 
% \\
% TAPE     & 10.8329 & 0.6106 & 1.2126 & 7.3417  & 1.2469 & 0.9717 
% % &  0.7019 
% \\
% ProtT5   & \textbf{3.1166}  & \textbf{0.1049} & 0.2275 & \textbf{1.8990}  & \textbf{0.2466} & 0.2045 
% % & 0.8803 
% \\
% \bottomrule
% \end{tabular}
% }
% \caption{Hyperbolicity and Ultrametricity values (max, average, and standard deviation) for Breast Cancer dataset.}
% \label{tab_breastcancer_metrics}
% \end{table}


% The hyperbolicity results for the PDB186 dataset, as presented in Table~\ref{tab_pdb186_hyperbolicity_ultramtricity_CCC}, indicate an average $\delta$ value of $0.4178$ with a standard deviation of $0.8754$, which is slightly higher than the synthetic dataset (note that a lower $\delta$ value is better). 
% This larger hyperbolicity suggests that the underlying structure of the Seqvec embeddings deviates more significantly from a tree-like structure, which could be due to the complex nature of DNA-protein binding interactions and the way Seqvec computed the embeddings. Whereas ProtT5 embeddings exhibit a tree-like structure, as evident from hyperbolicity and ultrametricity values. The higher $\delta$ value reflects the intricate and possibly non-uniform relationships within the dataset, where the distances between certain points do not conform to the idealized tree metric as closely as in the synthetic data. Additionally, the substantial standard deviation highlights the variability in the hyperbolicity across different parts of the dataset, suggesting that some regions may be more tree-like while others exhibit more complex, non-metric space characteristics. This complexity is expected in biological data, where the underlying interactions and dependencies are often non-linear and can vary significantly across different sequences. 


% The relationship between Euclidean and Poincaré distance matrices further illuminates our results. Embeddings computed using Euclidean distances show substantially higher NJ scores than their Poincaré counterparts, yet the relative ordering remains consistent: ProtT5 maintains the lowest values in both spaces, while SeqVec shows the highest. This suggests that while the absolute scale differs between distance metrics, the underlying hierarchical organization captured by each embedding method remains stable.



% \textcolor{green}{
% For PDB186 dataset, across all metrics, ProtT5 and TAPE exhibit consistently stronger hierarchical structure, characterized by stable negative Ricci curvature, low persistence lifetimes, and moderate intrinsic dimensionality, while SeqVec shows excessive multiscale fragmentation, and ESM exhibits numerical instability in fractal estimates.}

% \textcolor{green}{
% On the Lung Cancer dataset, ProtT5 again exhibits the strongest hierarchical organization, characterized by minimal persistence lifetimes and compact intrinsic dimensionality. TAPE consistently demonstrates low fractal dimensionality and stable curvature, while ESM shows stronger negative curvature but substantially higher intrinsic complexity. SeqVec exhibits weaker hierarchical consolidation with large persistence lifetimes and higher dimensional spread.
% }

% \textcolor{green}{
% On the Breast Cancer dataset, ProtT5 again exhibits the strongest hierarchical structure, with the smallest persistence lifetimes and compact intrinsic dimensionality. TAPE consistently demonstrates low fractal dimension and stable curvature, indicating well-formed multiscale clustering. ESM shows stronger negative curvature but higher intrinsic complexity, while SeqVec exhibits weaker hierarchical consolidation.}

\subsection{Classification Results}

% \textcolor{red}{
% Relative square deviation (Normalized)
% Root Mean Square Deviation
% Minimum evolution
% IQR
% }

The classification results on the PDB186 dataset are shown in Table~\ref{tbl_all_avg_data_results_pdb186}. ProtT5 achieves the highest ROC-AUC (0.7968 with LR), followed by ESM2 (0.7453), while SeqVec shows poor performance (0.6173). We can see that across other metrics as well, ProtT5 consistently outperforms SeqVec.
% This superior performance corresponds directly to their geometric properties: ProtT5's $\delta$-hyperbolicity and ultrametricity values are substantially lower than those of SeqVec (3.2018 and 16.6730), which achieves only 0.6173 ROC-AUC. The elevated hyperbolicity and ultrametricity of SeqVec indicate significant deviation from tree-like structures, resulting in embeddings that fail to organize DNA-protein binding relationships in a hierarchically meaningful way.
% }
% The highest ROC-AUC is observed for the LR model for ProtT5 embeddings at $0.7968$, followed by ESM2 embeddings for the LR model at $0.7453$. These models also perform well in terms of precision, recall, and F1 scores, indicating their ability to effectively classify the data. On the other hand, the seqvec embeddings show the lowest performance across all metrics, with an ROC-AUC of $0.6173$, which is consistent with the challenges of capturing complex relationships within the data using seqvec as justified by the $\delta$-hyperbolicity and ultrametricity values for the PDB186 dataset.
% The relatively high standard deviation observed in models like MLP (e.g., 0.04188 for accuracy) suggests variability in performance, which could be linked to the dataset's inherent complexity and non-linearity, as indicated by the higher hyperbolicity value ($\delta$) observed for PDB186.
\begin{table}[h!]
      \centering
      \scriptsize
      \resizebox{0.9\textwidth}{!}{
\begin{tabular}{@{\extracolsep{6pt}}p{0.5cm}lp{0.5cm}p{0.5cm}p{0.6cm}p{0.7cm}p{0.7cm}p{0.5cm}p{1.5cm}}
    \toprule
        \multirow{2}{*}{Embed.} & \multirow{2}{*}{Algo.} & \multirow{2}{*}{Acc. $\uparrow$} & \multirow{2}{*}{Prec. $\uparrow$} & \multirow{2}{*}{Recall $\uparrow$} & \multirow{2}{1.2cm}{F1 (Weig.) $\uparrow$} & \multirow{2}{1.2cm}{F1 (Macro) $\uparrow$} & \multirow{2}{1cm}{ROC AUC $\uparrow$} & Train Time (sec.) $\downarrow$\\
        \midrule \midrule
        % \cmidrule{2-9}  
      \multirow{7}{1.2cm}{SeqVec}
& SVM & 0.5786 & 0.6045 & 0.5786 & 0.5720 & 0.5755 & 0.5917 & 1692.2908		\\
& NB & 0.5036 & 0.5538 & 0.5036 & 0.3878 & 0.3836 & 0.5014 & \underline{132.7588}       \\
& MLP & 0.5607 & 0.5934 & 0.5607 & 0.5088 & 0.4911 & 0.5330 & 845.1396      \\
& KNN & 0.5107 & 0.3741 & 0.5107 & 0.3815 & 0.3812 & 0.5135 & 1808.8480     \\
& RF & 0.5857 & 0.5961 & 0.5857 & 0.5857 & 0.5856 & 0.5908 & 100.2860       \\
& LR & \underline{0.6179} & \underline{0.6216} & \underline{0.6179} & \underline{0.6184} & \underline{0.6159} & \underline{0.6173} & 1024.6331      \\
& DT & 0.5536 & 0.5550 & 0.5536 & 0.5531 & 0.5494 & 0.5507 & 277.4930       \\
 
\midrule
\multirow{7}{1.2cm}{ESM2}
& SVM & 0.7321 & 0.7520 & 0.7321 & 0.7327 & 0.7319 & 0.7420 & 0.0138	\\
& NB & 0.6321 & 0.6359 & 0.6321 & 0.6325 & 0.6273 & 0.6289 & \underline{0.0034}     \\
& MLP & 0.7107 & 0.7274 & 0.7107 & 0.7116 & 0.7103 & 0.7190 & 0.7072    \\
& KNN & 0.6214 & 0.6531 & 0.6214 & 0.6159 & 0.6178 & 0.6354 & 0.0171    \\
& RF & 0.6964 & 0.7231 & 0.6964 & 0.6952 & 0.6951 & 0.7080 & 0.3342     \\
& LR & \underline{0.7357} & \underline{0.7561} & \underline{0.7357} & \underline{0.7361} & \underline{0.7353} & \underline{0.7453} & 0.0452     \\
& DT & 0.5857 & 0.5986 & 0.5857 & 0.5851 & 0.5821 & 0.5894 & 0.0674     \\

 
\midrule
\multirow{7}{1.2cm}{TAPE}
& SVM & 0.6286 & 0.6411 & 0.6286 & 0.6301 & 0.6269 & 0.6317 & 0.0185	\\
& NB & 0.6893 & 0.6965 & 0.6893 & 0.6897 & 0.6858 & 0.6896 & \underline{\textbf{0.0031}}     \\
& MLP & \underline{0.6964} & \underline{0.7170} & \underline{0.6964} & \underline{0.6950} & \underline{0.6939} & \underline{0.7032} & 33.7174   \\
& KNN & 0.6000 & 0.6134 & 0.6000 & 0.5991 & 0.5954 & 0.6023 & 0.5691    \\
& RF & 0.6714 & 0.6974 & 0.6714 & 0.6696 & 0.6684 & 0.6795 & 0.4725     \\
& LR & 0.6286 & 0.6398 & 0.6286 & 0.6301 & 0.6268 & 0.6315 & 0.1313     \\
& DT & 0.5179 & 0.5372 & 0.5179 & 0.5153 & 0.5149 & 0.5259 & 0.0619     \\

\midrule

\multirow{7}{1.2cm}{ProtT5}
& SVM & 0.7679 & 0.7958 & 0.7679 & 0.7681 & 0.7675 & 0.7813 & 0.0235	\\
& NB & 0.7250 & 0.7385 & 0.7250 & 0.7228 & 0.7181 & 0.7261 & \underline{0.0035}     \\
& MLP & 0.7464 & 0.7558 & 0.7464 & 0.7459 & 0.7415 & 0.7465 & 38.2089   \\
& KNN & 0.6821 & 0.6928 & 0.6821 & 0.6818 & 0.6784 & 0.6843 & 0.6559    \\
& RF & 0.7036 & 0.7291 & 0.7036 & 0.7029 & 0.7015 & 0.7136 & 0.3849     \\
& LR & \underline{\textbf{0.7893}} & \underline{\textbf{0.8029}} & \underline{\textbf{0.7893}} & \underline{\textbf{0.7903}} & \underline{\textbf{0.7884}} & \underline{\textbf{0.7968}} & 0.1281     \\
& DT & 0.5607 & 0.5826 & 0.5607 & 0.5585 & 0.5567 & 0.5687 & 0.0841     \\
     
    \bottomrule
  \end{tabular}  
    }
    \caption{Average classification results for different models and algorithms for \textbf{PDB186} dataset.}
    \label{tbl_all_avg_data_results_pdb186}
\end{table}
Table~\ref{tbl_all_avg_data_results_LC} shows the classification results for the Lung Cancer dataset. Here, ESM2 achieves the highest ROC-AUC (0.8568 with SVM), followed by ProtT5 (0.8311 with LR). SeqVec and TAPE consistently show the lowest performance, with only marginal differences between them across other metrics.
% ESM2's moderate $\delta$-hyperbolicity (1.0897) and ultrametricity (1.4986) values position it between the highly tree-like ProtT5 and the poorly structured SeqVec (0.5904 and 1.2987). The consistently lower performance of SeqVec (0.8212 ROC-AUC) across all classifiers reinforces the association between greater geometric deviation and reduced classification performance.
% }
% The highest ROC-AUC is observed for the SVM model using ESM2 embeddings at $0.8568$, followed closely by the RF model using ProtT5 embeddings at $0.8311$. These models also perform well in terms of precision, recall, and F1 scores, indicating their ability to effectively classify the data. On the other hand, the SeqVec embeddings show the lowest performance across all metrics, with the worst performing ROC-AUC of $0.8212$ for SeqVec embeddings, which is consistent with the challenges of capturing complex relationships within the data using SeqVec as justified by the $\delta$-hyperbolicity and ultrametricity values for Lung Cancer datasets.
% The relatively high standard deviation observed in models like MLP (e.g., 0.04188 for accuracy) suggests variability in performance, which could be linked to the dataset's inherent complexity and non-linearity, as indicated by the higher hyperbolicity value ($\delta$) observed for the Lung cancer dataset.
\begin{table}[h!]
      \centering
      \scriptsize
      \resizebox{0.9\textwidth}{!}{
\begin{tabular}{@{\extracolsep{6pt}}p{0.5cm}lp{0.5cm}p{0.5cm}p{0.6cm}p{0.7cm}p{0.7cm}p{0.5cm}p{1.5cm}}
    \toprule
        \multirow{2}{*}{Embed.} & \multirow{2}{*}{Algo.} & \multirow{2}{*}{Acc. $\uparrow$} & \multirow{2}{*}{Prec. $\uparrow$} & \multirow{2}{*}{Recall $\uparrow$} & \multirow{2}{1.2cm}{F1 (Weig.) $\uparrow$} & \multirow{2}{1.2cm}{F1 (Macro) $\uparrow$} & \multirow{2}{1cm}{ROC AUC $\uparrow$} & Train Time (sec.) $\downarrow$\\
        \midrule \midrule
        % \cmidrule{2-9}  

      \multirow{7}{1.2cm}{SeqVec}
& SVM & \underline{0.9240} & \underline{0.9212} & \underline{0.9240} & \underline{0.9202} & \underline{0.6836} & \underline{0.8212} & 29.3541		\\
& NB & 0.6347 & 0.7938 & 0.6347 & 0.6212 & 0.2585 & 0.6396 & \underline{3.6623}         \\
& MLP & 0.9048 & 0.9101 & 0.9048 & 0.9056 & 0.6513 & 0.8121 & 142.6836      \\
& KNN & 0.8812 & 0.8660 & 0.8812 & 0.8683 & 0.5436 & 0.7165 & 9.0200        \\
& RF & 0.9085 & 0.8955 & 0.9085 & 0.8977 & 0.6329 & 0.7649 & 7.2720         \\
& LR & 0.7882 & 0.8925 & 0.7882 & 0.8254 & 0.5433 & 0.8066 & 1005.2145      \\
& DT & 0.8531 & 0.8692 & 0.8531 & 0.8593 & 0.5348 & 0.7469 & 23.7610        \\

\midrule
\multirow{7}{1.2cm}{ESM2}
& SVM & \underline{\textbf{0.9506}} & \underline{\textbf{0.9474}} & \underline{\textbf{0.9506}} & \underline{\textbf{0.9475}} & \underline{\textbf{0.7335}} & \underline{\textbf{0.8568}} & 11.9609		\\
& NB & 0.6472 & 0.8906 & 0.6472 & 0.7109 & 0.5231 & 0.8066 & \underline{0.6002}         \\
& MLP & 0.9365 & 0.9453 & 0.9365 & 0.9389 & 0.7428 & 0.8698 & 9.4761        \\
& KNN & 0.9122 & 0.9259 & 0.9122 & 0.9144 & 0.6687 & 0.8300 & 1.2228        \\
& RF & 0.9380 & 0.9352 & 0.9380 & 0.9342 & 0.7144 & 0.8296 & 2.8020         \\
& LR & 0.7823 & 0.8957 & 0.7823 & 0.8171 & 0.5472 & 0.8243 & 156.7906       \\
& DT & 0.9011 & 0.9152 & 0.9011 & 0.9039 & 0.6469 & 0.8074 & 5.3667         \\

\midrule
\multirow{7}{1.2cm}{TAPE}
& SVM & 0.9166 & 0.9181 & 0.9166 & 0.9165 & 0.6177 & 0.7967 & 0.0394	\\
& NB & 0.9018 & 0.9204 & 0.9018 & 0.9082 & 0.6461 & \underline{0.8306} & \underline{\textbf{0.0095}}     \\
& MLP & 0.9085 & 0.9152 & 0.9085 & 0.9102 & 0.6194 & 0.7943 & 1.7274    \\
& KNN & 0.9122 & 0.9143 & 0.9122 & 0.9107 & 0.5994 & 0.7771 & 0.0265    \\
& RF & 0.9144 & 0.9065 & 0.9144 & 0.9054 & 0.6336 & 0.7546 & 0.9389     \\
& LR & \underline{0.9173} & \underline{0.9252} & \underline{0.9173} & \underline{0.9200} & \underline{0.6564} & 0.8221 & 0.7910     \\
& DT & 0.8657 & 0.8873 & 0.8657 & 0.8748 & 0.5586 & 0.7576 & 0.3775     \\

\midrule
\multirow{7}{1.2cm}{ProtT5}
& SVM & 0.9277 & 0.9313 & 0.9277 & 0.9277 & 0.6664 & 0.8290 & 0.0513	\\
& NB & 0.9151 & 0.9226 & 0.9151 & 0.9175 & 0.6461 & 0.8232 & \underline{0.0125}     \\
& MLP & 0.9196 & 0.9263 & 0.9196 & 0.9219 & 0.6375 & 0.8155 & 3.5043    \\
& KNN & 0.9188 & 0.9169 & 0.9188 & 0.9165 & 0.5826 & 0.7817 & 0.0649    \\
& RF & \underline{0.9321} & 0.9293 & \underline{0.9321} & 0.9258 & 0.6520 & 0.7997 & 1.0917     \\
& LR & 0.9277 & \underline{0.9326} & 0.9277 & \underline{0.9288} & \underline{0.6724} & \underline{0.8311} & 1.0877     \\
& DT & 0.8841 & 0.8977 & 0.8841 & 0.8887 & 0.5800 & 0.7800 & 0.4553     \\
    \bottomrule
  \end{tabular}  
    }
    \caption{Average classification results for different models and algorithms for \textbf{Lung Cancer} dataset.}
    \label{tbl_all_avg_data_results_LC}
\end{table}
The classification results for the Breast Cancer dataset are shown in Table~\ref{tbl_all_avg_data_results_BC}, and they reinforce the trend observed previously. ESM2 attains the highest ROC-AUC (0.8385), followed by ProtT5 at 0.8195, while SeqVec again shows the weakest performance (0.7597). Across all other metrics, SeqVec consistently lags behind by a noticeable margin, further highlighting the superiority of transformer-based embeddings for classification tasks.

% \textcolor{green}{
% The Breast Cancer dataset results in Table~\ref{tbl_all_avg_data_results_BC} further validate our observations, with ESM2 achieving 0.8385 ROC-AUC followed by ProtT5 reaching 0.8195. SeqVec again demonstrates the weakest performance at 0.7597, corresponding to its higher $\delta$-hyperbolicity (0.5804) and ultrametricity (1.3425) in Euclidean distance compared to ProtT5. 
% The ESM2 not only exhibits a clear tree-like structure in the Poincaré distance matrix for both the Lung and Breast cancer datasets, but this hierarchical organization is also reflected in their strong classification performance.
% Also, the Neighbor Joining scores corroborate these findings: ProtT5 consistently shows lower NJ values (318.63 for PDB186, 2296.54 for Lung Cancer, 2487.96 for Breast Cancer) compared to SeqVec (7276.80, 28674.59, 30408.33), indicating stronger adherence to additive tree structure.
% }

% Table~\ref{tbl_all_avg_data_results_BC} shows classification results for the Breast Cancer dataset. The highest ROC-AUC is observed for the SVM model for ESM2 embeddings at $0.8385$, followed closely by ProtT5 embeddings for the RF model at $0.8195$. These models also perform well in terms of precision, recall, and F1 scores, indicating their ability to effectively classify the data. On the other hand, similar to Lung Cancer and PDB186, the Seqvec embeddings show the lowest performance across all metrics, with the worst performing ROC-AUC of $0.7597$, which is consistent with the challenges of capturing complex relationships within the data using seqvec as justified by the $\delta$-hyperbolicity and ultrametricity values for Lung Cancer datasets.

\begin{table}[h!]
      \centering
      \scriptsize
      \resizebox{0.9\textwidth}{!}{
\begin{tabular}{@{\extracolsep{6pt}}p{0.5cm}lp{0.5cm}p{0.5cm}p{0.6cm}p{0.7cm}p{0.7cm}p{0.5cm}p{1.5cm}}
    \toprule
        \multirow{2}{*}{Embed.} & \multirow{2}{*}{Algo.} & \multirow{2}{*}{Acc. $\uparrow$} & \multirow{2}{*}{Prec. $\uparrow$} & \multirow{2}{*}{Recall $\uparrow$} & \multirow{2}{1.2cm}{F1 (Weig.) $\uparrow$} & \multirow{2}{1.2cm}{F1 (Macro) $\uparrow$} & \multirow{2}{1cm}{ROC AUC $\uparrow$} & Train Time (sec.) $\downarrow$\\
        \midrule \midrule
        % \cmidrule{2-9}  
      \multirow{7}{1.2cm}{SeqVec}
& SVM & \underline{0.9004} & \underline{0.8995} & \underline{0.9004} & \underline{0.8977} & \underline{0.5682} & \underline{0.7597} & 1724.5	\\
& NB & 0.4996 & 0.7564 & 0.4996 & 0.5289 & 0.2087 & 0.5882 & \underline{4.2771} \\
& MLP & 0.8716 & 0.8833 & 0.8716 & 0.8740 & 0.5218 & 0.7447 & 277.84 \\
& KNN & 0.7761 & 0.8375 & 0.7761 & 0.7969 & 0.4523 & 0.6859 & 175.88 \\
& RF & 0.8814 & 0.8649 & 0.8814 & 0.8671 & 0.5349 & 0.7082 & 11.531 \\
& LR & 0.7719 & 0.8811 & 0.7719 & 0.8164 & 0.5014 & 0.7528 & 1836.9 \\
& DT & 0.8133 & 0.8245 & 0.8133 & 0.8179 & 0.4515 & 0.6745 & 37.130 \\

\midrule
\multirow{7}{1.2cm}{ESM2} 
& SVM & \underline{\textbf{0.9333}} & \underline{\textbf{0.9335}} & \underline{\textbf{0.9333}} & \underline{\textbf{0.9319}} & \underline{\textbf{0.6843}} & \underline{\textbf{0.8385}} & 8.0102 \\
& NB & 0.6681 & 0.8646 & 0.6681 & 0.7150 & 0.5013 & 0.7977 & \underline{0.7043} \\
& MLP & 0.9067 & 0.9182 & 0.9067 & 0.9105 & 0.6353 & 0.8032 & 12.827 \\
& KNN & 0.9060 & 0.9089 & 0.9060 & 0.9047 & 0.6380 & 0.8268 & 1.1774 \\
& RF & 0.9326 & 0.9257 & 0.9326 & 0.9261 & 0.6589 & 0.7914 & 2.7038 \\
& LR & 0.7993 & 0.8796 & 0.7993 & 0.8239 & 0.5591 & 0.8177 & 154.53 \\
& DT & 0.8877 & 0.8889 & 0.8877 & 0.8871 & 0.6237 & 0.7948 & 3.7809 \\

\midrule
\multirow{7}{1.2cm}{TAPE}
& SVM & 0.8982 & 0.8968 & 0.8982 & 0.8962 & 0.5964 & 0.7813 & 0.0930	\\
& NB & 0.8575 & 0.8852 & 0.8575 & 0.8691 & 0.5760 & 0.7887 & \underline{\textbf{0.0122}}     \\
& MLP & 0.8884 & 0.8943 & 0.8884 & 0.8901 & 0.5825 & 0.7815 & 284.43  \\
& KNN & 0.8821 & 0.8716 & 0.8821 & 0.8745 & 0.5119 & 0.7345 & 1.1070    \\
& RF & 0.8870 & 0.8683 & 0.8870 & 0.8738 & 0.5615 & 0.7374 & 2.2642     \\
& LR & \underline{0.9046} & \underline{0.9082} & \underline{0.9046} & \underline{0.9030} & \underline{0.6161} & \underline{0.7961} & 1.9700     \\
& DT & 0.8386 & 0.8511 & 0.8386 & 0.8433 & 0.5123 & 0.7262 & 0.9652     \\


\midrule
\multirow{7}{1.2cm}{ProtT5}
& SVM & 0.8975 & 0.9011 & 0.8975 & 0.8973 & 0.6425 & 0.8086 & 0.0743	\\
& NB & 0.8842 & 0.8902 & 0.8842 & 0.8865 & 0.5744 & 0.7770 & \underline{0.0142}     \\
& MLP & 0.8961 & \underline{0.9015} & 0.8961 & \underline{0.8977} & \underline{0.6509} & \underline{0.8195} & 3.3378    \\
& KNN & 0.8933 & 0.8796 & 0.8933 & 0.8855 & 0.5230 & 0.7425 & 0.0496    \\
& RF & \underline{0.9018} & 0.8892 & \underline{0.9018} & 0.8934 & 0.5996 & 0.7651 & 1.3799     \\
& LR & 0.8947 & 0.8973 & 0.8947 & 0.8951 & 0.6131 & 0.8004 & 1.0668     \\
& DT & 0.8484 & 0.8535 & 0.8484 & 0.8497 & 0.5021 & 0.7278 & 0.6561     \\
    \bottomrule
  \end{tabular}  
    }
    \caption{Average classification results for different models and algorithms for \textbf{Breast Cancer} dataset.}
    \label{tbl_all_avg_data_results_BC}
\end{table}


% \textcolor{blue}{
% The classification results demonstrate a strong correlation between geometric properties of embeddings and their performance on downstream tasks. Across all three datasets, ProtT5 embeddings consistently achieve superior classification performance, which aligns with their geometric characteristics: lowest average $\delta$-hyperbolicity values (0.0418 for PDB186, 0.1010 for Lung Cancer, 0.1073 for Breast Cancer) and ultrametricity values (0.1300, 0.2413, 0.2458, respectively, for Euclidean distance matrices). While the lowest average $\delta$-hyperbolicity values (0.0030 for PDB186, 0.0016 for Lung Cancer, 0.0017 for Breast Cancer) and ultrametricity values (0.0083, 0.0022, 0.0023, respectively, for Poincaré distance matrices). These low values indicate that ProtT5 embeddings and ESM2 embeddings in respective distance matrices exhibit strong tree-like structures that effectively capture hierarchical biological relationships inherent in protein sequences.
% }

% \textcolor{green}{
% The relationship between Euclidean and Poincaré distance matrices further illuminates our results. Embeddings computed using Euclidean distances show substantially higher NJ scores than their Poincaré counterparts, yet the relative ordering remains consistent: ProtT5 maintains the lowest values in both spaces, while SeqVec shows the highest. This suggests that while the absolute scale differs between distance metrics, the underlying hierarchical organization captured by each embedding method remains stable and predictive of classification performance.
% }

% \textcolor{green}{
% The consistent alignment between geometric properties ($\delta$-hyperbolicity, ultrametricity, NJ scores) and classification performance across all datasets provides strong empirical evidence that tree-likeness in embedding spaces is beneficial for capturing the hierarchical nature of protein sequence relationships. Models that generate more tree-like embeddings create representations where biologically related sequences are organized in a manner that facilitates accurate classification by downstream machine learning algorithms.
% }


\subsection{Clustering Results}
Table~\ref{tbl_clustering_PDB_186} presents the clustering results for the PDB186 dataset, while Table~\ref{tbl_clustering_Lung_Cancer} and Table~\ref{tbl_clustering_Breast_Cancer} show the results for the Lung Cancer and Breast Cancer datasets, respectively.

For the PDB186 dataset, SeqVec shows an interesting contrast: it performs best on clustering metrics such as the Silhouette Coefficient and Davies–Bouldin Score in all clustering algorithms, yet it is the weakest model in the classification setting. ESM2, which ranked second in classification, shows more consistent behavior, particularly on the Calinski–Harabasz Score, where it performs well for Agglomerative and k-means clustering. This contrast highlights how models that excel in supervised classification do not necessarily produce the most cluster-friendly embedding spaces, and vice versa.

% For the Lung Cancer dataset, ProtT5 embeddings consistently achieve the strongest performance across all clustering metrics—Silhouette Coefficient, Calinski–Harabasz Score, and Davies–Bouldin Score—and across both Agglomerative and k-means clustering. In contrast, results for the k-modes algorithm align more closely with the classification setting: ESM2, which performed best in supervised tasks (though only by a marginal margin over ProtT5), also achieves stronger k-modes scores, particularly on the Silhouette Coefficient and Calinski–Harabasz Score. This suggests that while ProtT5 produces highly clusterable representations in continuous spaces, ESM2 better captures structure relevant to categorical clustering, reflecting the small but meaningful differences observed in classification.

For the Lung Cancer dataset, ProtT5 performs better in all clustering metrics for Agglomerative and k-means. However, the k-modes findings are more in line with classification: ESM2, which is marginally more effective in supervised tasks, also receives higher k-modes scores, particularly on the Silhouette and Calinski–Harabasz scores. We can say that ProtT5 yields more cluster-friendly continuous embeddings, whereas ESM2 captures a structure better suited to categorical clustering.

% This divergence highlights an important insight where a strong tree-like structure in the embedding space is beneficial for complex or hierarchical datasets, but may not be necessary for simpler binary classification tasks.

We find a trend comparable to the Lung Cancer results in the Breast Cancer dataset. Although ESM2 outperforms ProtT5 in classification, it performs better in k-modes, particularly on the Silhouette and Calinski–Harabasz scores. ProtT5 has the best clustering performance for Agglomerative and k-means. This again exhibits a similar pattern: ProtT5 produces more cluster-friendly continuous embeddings, while ESM2 captures structure more closely aligned with categorical clustering and supervised classification tasks.

% We can see that the performance of ProtT5 embeddings varies across datasets in relation to both clustering quality metrics and structural properties of the embedding space. On the Lung Cancer and Breast Cancer datasets, ProtT5 embeddings consistently yield the best results across all clustering evaluation metrics—Silhouette Coefficient, Calinski-Harabasz Score, and Davies-Bouldin Score. ProtT5 consistently performs the best for Agglomerative and k-means clustering algorithms. Moreover, these embeddings also exhibit the highest values for Ultrametricity and Hyperbolicity, indicating that the learned representations are hierarchically organized and tree-like in structure. Although tree-like embeddings are not universally better, they are very useful for tasks involving hierarchy. They do not ensure improved clustering performance unless that hierarchical structure is reflected in the clustering task.
% This can be seen with the PDB186 dataset; here, ProtT5 does not achieve the best clustering performance, the SeqVec embeddings that perform best according to clustering metrics exhibit the worst hyperbolicity and ultrametricity values. This divergence highlights an important insight where a strong tree-like structure in the embedding space is beneficial for complex or hierarchical datasets, but may not be necessary for simpler binary classification tasks.

% The observed contradiction arises because clustering metrics evaluate local separation, while hyperbolicity and ultrametricity reflect global geometric properties, particularly whether the space resembles a tree or hierarchy. In simple binary tasks like PDB186, strong cluster separation does not require a tree-like embedding space. Clustering metrics like Silhouette Score measure how well separated the clusters are locally. We can have well-separated clusters without the embedding space being tree-like. Moreover, PDB186 is a binary classification problem with relatively clear separation between the two classes. Such problems may benefit more from flatter, linearly separable embeddings than from hierarchical structures. In contrast, Lung and Breast Cancer datasets may involve multi-class or more nuanced biological variation, where tree-like relationships help capture underlying semantics or phenotype similarities.


\begin{table}[h!]
  \centering
  \scriptsize
  \resizebox{0.9\textwidth}{!}{
  \begin{tabular}{p{1.8cm}p{1.2cm}p{1.2cm}p{1.4cm}p{1.4cm}p{1.3cm}p{1.8cm}}
    \toprule
    & & \multicolumn{3}{c}{Evaluation Metrics} \\
    \cmidrule{3-5}
    Algorithm & Embedding & Silhouette Coefficient $\uparrow$ & Calinski-Harabasz Score $\uparrow$ & Davies-Bouldin Score $\downarrow$ & Clustering Runtime $\downarrow$\\
    \midrule \midrule
    \multirow{4}{*}{Agglomerative} & 
     Seqvec & \textbf{\underline{0.603}} & 9.872 & \textbf{\underline{0.278}} &  \textbf{\underline{0.01}} \\
    & ESM2 & 0.230 & \underline{44.368} & 1.348 &  0.02 \\
    & TAPE & 0.136 & 12.068 & 3.665 & \textbf{\underline{0.01}} \\
    & ProtT5 & 0.117 & 10.429 &  2.306 & 0.03 \\
    \midrule
    \multirow{4}{*}{$k$-means} & 
     Seqvec &  \underline{0.384} & 3.407 & \underline{0.469} &  1.46 \\
    & ESM2 & 0.178 & \textbf{\underline{51.234}} & 1.820 &  1.46 \\
    & TAPE & 0.096 & 13.810 & 3.567 & \underline{1.20} \\
    & ProtT5 & 0.055 & 11.117 &  3.963 & 1.78 \\
    \midrule
    \multirow{4}{*}{$k$-modes} & 
     Seqvec & \underline{0.312} & \underline{2.545} & \underline{0.543} &  \underline{0.73} \\
    & ESM2 & 0.161 & 2.118 & 0.657 &  3.58 \\
    & TAPE & 0.033 & 1.167 & 0.895 &  2.32 \\
    & ProtT5 & -0.075 & 0.698 & 1.161 &  2.73 \\
    \bottomrule
  \end{tabular}
  }
  \caption{Internal clustering quality metrics for $k$-means
    and $k$-modes on different embeddings on the
    PDB186. Best values are shown in bold.}
  \label{tbl_clustering_PDB_186}
\end{table}

\begin{table}[h!]
  \centering
    \scriptsize
  \resizebox{0.9\textwidth}{!}{
  \begin{tabular}{p{1.8cm}p{1.2cm}p{1.2cm}p{1.4cm}p{1.4cm}p{1.3cm}}
    \toprule
    & & \multicolumn{3}{c}{Evaluation Metrics} \\
    \cmidrule{3-5}
    Algorithm & Embedding & Silhouette Coefficient $\uparrow$ & Calinski-Harabasz Score $\uparrow$ & Davies-Bouldin Score $\downarrow$ & Clustering Runtime $\downarrow$\\
    \midrule \midrule
     \multirow{4}{*}{Agglomerative} & 
     Seqvec & -0.002 & 36.713 & 3.792 & \textbf{\underline{0.22}} \\
    & ESM2 & 0.048 & 75.235 & 3.012 &  9.67 \\
    & TAPE & 0.082 & 80.764 & 2.709 &  0.23 \\
    & ProtT5 & \underline{0.177} & \underline{88.294} & \underline{2.030} &  0.25 \\
    \midrule
    \multirow{4}{*}{$k$-means} & 
     Seqvec & 0.026 & 38.255 & 3.922 & 2.29 \\
    & ESM2 & 0.070 & 79.774 & 3.370 &  5.95 \\
    & TAPE & 0.118 & \textbf{\underline{98.614}} & 2.527 &  2.54 \\
    & ProtT5 & \textbf{\underline{0.197}} & 94.129 & \underline{1.947} &  \underline{1.58} \\
    \midrule
    \multirow{3}{*}{$k$-modes} & 
     Seqvec & -0.038 & 0.947 & \underline{1.023} &  \underline{16.12} \\
    & ESM2 & \underline{0.035} & \underline{44.883} & 6.223 &  1642.4 \\
    & TAPE &  -0.168 & 0.837 & 1.087 &  18.40 \\
    & ProtT5 & -0.162 & 1.388 & \textbf{\underline{1.023}} & 22.34 \\
    \bottomrule
  \end{tabular}
  }
  \caption{Internal clustering quality metrics for $k$-means
    and $k$-modes on different embeddings on the
    Lung Cancer. Best values are shown in bold.}
  \label{tbl_clustering_Lung_Cancer}
\end{table}

\begin{table}[h!]
  \centering
    \scriptsize
  \resizebox{0.9\textwidth}{!}{
 \begin{tabular}{p{1.8cm}p{1.2cm}p{1.2cm}p{1.4cm}p{1.4cm}p{1.3cm}}
    \toprule
    & & \multicolumn{3}{c}{Evaluation Metrics} \\
    \cmidrule{3-5}
    Algorithm & Embedding & Silhouette Coefficient $\uparrow$ & Calinski-Harabasz Score $\uparrow$ & Davies-Bouldin Score $\downarrow$ & Clustering Runtime $\downarrow$\\
    \midrule	\midrule	
     \multirow{4}{*}{Agglomerative} & 
     Seqvec & -0.002 & 38.110 & 3.868 & 0.30 \\
    & ESM2 & 0.048 & 78.754 & 3.064 & 10.75 \\
    & TAPE & 0.096 & 87.078 & 2.826 & 0.17 \\
    & ProtT5 & \textbf{\underline{0.206}} & \underline{100.927} & \underline{2.003} & \textbf{\underline{0.14}} \\
    \midrule
    \multirow{4}{*}{$k$-means} & 
     Seqvec & -0.004 & 37.709 & 3.765 &  2.41 \\
    & ESM2 & 0.072 & 84.914 & 3.380 &  5.71 \\
    & TAPE & 0.122 & \textbf{\underline{109.023}} & 2.605 & \underline{1.56} \\
    & ProtT5 & \underline{0.190} & 106.776 & \underline{2.021} &  2.32 \\
    \midrule
    \multirow{4}{*}{$k$-modes} & 
     Seqvec & -0.093 & 0.910 & 1.078 & 18.29 \\
    & ESM2 & \underline{0.037} & \underline{47.333} & 5.461 &  2277.34 \\
    & TAPE & -0.116 & 1.035 & 0.984 &  \underline{8.96} \\
    & ProtT5 & -0.169 & 2.356 & \textbf{\underline{0.957}} &  9.50 \\
    \bottomrule
  \end{tabular}
  }
  \caption{Internal clustering quality metrics for $k$-means
    and $k$-modes on different embeddings on the
    Breast Cancer. Best values are shown in bold.}
  \label{tbl_clustering_Breast_Cancer}
\end{table}


% \begin{table*}[ht!]
% \centering
% \tiny
% \setlength{\tabcolsep}{3pt}
% \renewcommand{\arraystretch}{0.95}

% % ===================== PDB186 =====================
% \begin{subtable}[t]{0.32\textwidth}
% \centering
% \caption{PDB186}
% \resizebox{\textwidth}{!}{
% \begin{tabular}{l l c c c c}
% \toprule
% Alg. & Emb. & Sil. $\uparrow$ & CH  $\uparrow$& DB  $\downarrow$& Time  $\downarrow$ \\
% \midrule
% \multirow{4}{*}{Agglo}
% & SeqVec & \textbf{\underline{0.603}} & 9.872 & \textbf{\underline{0.278}} & \textbf{\underline{0.010}} \\
% & ESM2   & 0.230 & \textbf{\underline{44.368}} & 1.348 & 0.020 \\
% & TAPE   & 0.136 & 12.068 & 3.665 & \textbf{\underline{0.010}} \\
% & ProtT5 & 0.117 & 10.429 & 2.306 & 0.030 \\
% \midrule
% \multirow{4}{*}{k-means}
% & SeqVec & \textbf{\underline{0.384}} & 3.407 & \textbf{\underline{0.469}} & 1.460 \\
% & ESM2   & 0.178 & \textbf{\underline{51.234}} & 1.820 & 1.460 \\
% & TAPE   & 0.096 & 13.810 & 3.567 & \textbf{\underline{1.200}} \\
% & ProtT5 & 0.055 & 11.117 & 3.963 & 1.780 \\
% \midrule
% \multirow{4}{*}{k-modes}
% & SeqVec & \textbf{\underline{0.312}} & \textbf{\underline{2.545}} & \textbf{\underline{0.543}} & \textbf{\underline{0.730}} \\
% & ESM2   & 0.161 & 2.118 & 0.657 & 3.580 \\
% & TAPE   & 0.033 & 1.167 & 0.895 & 2.320 \\
% & ProtT5 & -0.075 & 0.698 & 1.161 & 2.730 \\
% \bottomrule
% \end{tabular}}
% \end{subtable}
% \hfill
% % ===================== Lung Cancer =====================
% \begin{subtable}[t]{0.32\textwidth}
% \centering
% \caption{Lung Cancer}
% \resizebox{\textwidth}{!}{
% \begin{tabular}{l l c c c c}
% \toprule
% Alg. & Emb. & Sil. $\uparrow$ & CH  $\uparrow$& DB  $\downarrow$& Time  $\downarrow$ \\
% \midrule
% \multirow{4}{*}{Agglo}
% & SeqVec & -0.002 & 36.713 & 3.792 & \textbf{\underline{0.220}} \\
% & ESM2   & 0.048 & 75.235 & 3.012 & 9.670 \\
% & TAPE   & 0.082 & 80.764 & 2.709 & 0.230 \\
% & ProtT5 & \textbf{\underline{0.177}} & \textbf{\underline{88.294}} & \textbf{\underline{2.030}} & 0.250 \\
% \midrule
% \multirow{4}{*}{k-means}
% & SeqVec & 0.026 & 38.255 & 3.922 & 2.290 \\
% & ESM2   & 0.070 & 79.774 & 3.370 & 5.950 \\
% & TAPE   & 0.118 & \textbf{\underline{98.614}} & 2.527 & 2.540 \\
% & ProtT5 & \textbf{\underline{0.197}} & 94.129 & \textbf{\underline{1.947}} & \textbf{\underline{1.580}} \\
% \midrule
% \multirow{4}{*}{k-modes}
% & SeqVec & -0.038 & 0.947 & \textbf{\underline{1.023}} & \textbf{\underline{16.120}} \\
% & ESM2   & \textbf{\underline{0.035}} & \textbf{\underline{44.883}} & 6.223 & 1642.400 \\
% & TAPE   & -0.168 & 0.837 & 1.087 & 18.400 \\
% & ProtT5 & -0.162 & 1.388 & \textbf{\underline{1.023}} & 22.340 \\
% \bottomrule
% \end{tabular}}
% \end{subtable}
% \hfill
% % ===================== Breast Cancer =====================
% \begin{subtable}[t]{0.32\textwidth}
% \centering
% \caption{Breast Cancer}
% \resizebox{\textwidth}{!}{
% \begin{tabular}{l l c c c c}
% \toprule
% Alg. & Emb. & Sil. $\uparrow$ & CH  $\uparrow$& DB  $\downarrow$& Time  $\downarrow$ \\
% \midrule
% \multirow{4}{*}{Agglo}
% & SeqVec & -0.002 & 38.110 & 3.868 & 0.300 \\
% & ESM2   & 0.048 & 78.754 & 3.064 & 10.750 \\
% & TAPE   & 0.096 & 87.078 & 2.826 & 0.170 \\
% & ProtT5 & \textbf{\underline{0.206}} & \textbf{\underline{100.927}} & \textbf{\underline{2.003}} & \textbf{\underline{0.140}} \\
% \midrule
% \multirow{4}{*}{k-means}
% & SeqVec & -0.004 & 37.709 & 3.765 & 2.410 \\
% & ESM2   & 0.072 & 84.914 & 3.380 & 5.710 \\
% & TAPE   & 0.122 & \textbf{\underline{109.023}} & 2.605 & \textbf{\underline{1.560}} \\
% & ProtT5 & \textbf{\underline{0.190}} & 106.776 & \textbf{\underline{2.021}} & 2.320 \\
% \midrule
% \multirow{4}{*}{k-modes}
% & SeqVec & -0.093 & 0.910 & 1.078 & 18.290 \\
% & ESM2   & \textbf{\underline{0.037}} & \textbf{\underline{47.333}} & 5.461 & 2277.340 \\
% & TAPE   & -0.116 & 1.035 & 0.984 & \textbf{\underline{8.960}} \\
% & ProtT5 & -0.169 & 2.356 & \textbf{\underline{0.957}} & 9.500 \\
% \bottomrule
% \end{tabular}}
% \end{subtable}
% \caption{Clustering results for different embeddings across three datasets. Best values are shown in bold.}
% \label{tab:clustering_side_by_side_all}
% \end{table*}




% \subsection{Limitations}
% While this study provides valuable insights into the geometric properties of LLM-generated embeddings through the lens of $\delta$-hyperbolicity and ultrametricity, several limitations must be acknowledged. First, the evaluation of hyperbolicity and ultrametricity is highly dependent on the choice of metric and dataset, and the results may not generalize across different embedding types or application domains. This requires analyzing our method on a more diverse dataset, which we will do in the future.
% % Additionally, the computational complexity of computing hyperbolicity for large datasets can be significant, potentially limiting the scalability of this approach for more extensive health recommendation systems. 
% Furthermore, this work primarily focuses on a specific type of LLM and embedding method, which may not capture the full diversity of existing models. For this, we will explore additional LLMs from the literature that are fine-tuned on protein sequences.
% % Lastly, while we demonstrate the relationship between hyperbolicity and recommendation accuracy, the causality and underlying mechanisms warrant further investigation. 
% Future work should explore these aspects to provide a more comprehensive understanding of how geometric properties can be used in their Applications, such as health recommender systems or graph based tasks.


% However, studying the tree-like structure of graphs based solely on the value of the hyperbolicity may not be sufficient to characterize their closeness to a tree topology for two reasons. First, the hyperbolicity is a relative measure. For example, for a given graph $G=(V, E)$, a value of $δ(G)=10$ can be seen as too large when the structure of G can be fairly described as being far from a tree. However, when 
% the hyperbolicity $δ(G)=10$ looks much smaller, and indicates a tree-like structure. Second, small graph size and (or) small diameter directly yield low hyperbolicity; our classification should be sensitive to such cases. In other words, small δ(G) does not always suggest a graph with a tree-like structure; other graph attributes that might impact the hyperbolicity must be investigated. We find that the graph’s size size(G) and the graph’s diameter diam(G) play an important role in deciding how hyperbolic a given graph is.

\subsection{Overall Results - Observations}
Across all three datasets, we observe a pattern indicating the importance of geometry for downstream performance. The classification task aligns with a strong Poincaré hierarchy, in which ESM2 exhibits a better hyperbolic structure and consistently performs better on supervised tasks. Clustering aligns with Euclidean space, and ProtT5 exhibits the most separable Euclidean structure and gives better performance with Agglomerative and k-means clustering. Although it occasionally clusters well, SeqVec, the worst in classification, also has the most fragmented geometry, which explains its inconsistent behavior. Although TAPE exhibits a strong hyperbolic hierarchy, it does not translate into downstream performance, indicating that the hyperbolic structure alone is insufficient. In general, ESM2 excels at classification, ProtT5 at clustering, and SeqVec and TAPE exhibit weaker or less consistent geometric patterns.

% Overall, the PDB186 results show that geometric structure strongly shapes downstream performance. ProtT5, which achieves the highest classification ROC-AUC, also shows the strongest Euclidean tree-likeness, indicating that supervised accuracy is driven by good Euclidean separability. TAPE exhibits the most hierarchical Poincaré structure, but this does not translate into better task performance, showing that hyperbolic hierarchy alone is not predictive. SeqVec performs poorly in classification but well in clustering, reflecting its strong local curvature and high persistence entropy, which produce highly clusterable but less separable embeddings. Taken together, ProtT5 and TAPE display stronger hierarchical tendencies overall, while SeqVec shows fragmented multiscale geometry and ESM exhibits instability in fractal behavior.

% The Lung Cancer results show a clear alignment between each embedding's geometry and its downstream behavior. ESM2, which performs best in classification, also shows the cleanest hierarchical structure in Poincaré space, suggesting that its embeddings capture strong global organization useful for supervised learning. ProtT5, although slightly behind ESM2 in classification, becomes the most structured in Euclidean space—matching its superior performance in Agglomerative and k-means clustering. SeqVec and TAPE, which perform the weakest in classification, also show less coherent geometric structure. Overall, ESM2 aligns with classification strength, ProtT5 with cluster-friendly Euclidean organization, and the weaker models exhibit less consistent geometric patterns.

% The geometric patterns for the Breast Cancer dataset align with the differences seen in classification and clustering performance. ESM2, which performs the best in classification, also shows the most well-organized hierarchical structure in the Poincaré space, suggesting that its embeddings capture clearer global relationships useful for supervised learning. ProtT5, although slightly behind in classification, shows stronger structure in Euclidean space, matching its consistently strong performance in clustering methods like Agglomerative and k-means. SeqVec, the weakest model in classification, also appears the least organized geometrically, showing more fragmented and less coherent structure. Overall, ESM2’s global hierarchy supports its classification strength, ProtT5’s Euclidean organization supports its clustering strength, and SeqVec’s fragmented geometry reflects its weaker overall performance.

\section{Conclusion}
\label{sec_conclusion}
In this paper, we present a geometric framework for evaluating LLM embeddings using $\delta$-hyperbolicity, ultrametricity, Neighbor Joining, and complementary measures such as Ricci curvature, persistent homology, and fractal dimension. We found that the embedding geometry aligns with the downstream performance. We found that ESM2's hyperbolic structure supports strong classification performance, that ProtT5’s Euclidean separability explains its superior clustering, and that weaker models like SeqVec show fragmented geometric patterns consistent with lower task performance. The results highlight that no single embedding is universally optimal. Each model excels in different geometric regimes. Each model is better suited to a specific task, and this structural analysis can pave the way towards guided, informed designs and the selection of LLMs.
Our analysis of synthetic spaces further confirms how hyperbolicity varies across geometric settings, emphasizing the importance of understanding embedding structure. All things considered, geometric metrics provide a useful perspective for choosing and fine-tuning LLM embeddings. Future work can extend this framework to broader datasets and explore integrating geometric insights into model design and training.

% In this paper, we introduced an approach to evaluate the effectiveness of embeddings generated by a large language model (LLM) by assessing their $\delta$-hyperbolicity, ultrametricity, and Neighbor Joining. By leveraging these geometric measures, we aimed to uncover the underlying structure of LLM embeddings. Our findings indicate that the embeddings exhibit varying degrees of hyperbolicity, ultrametricity,  and Neighbor Joining (NJ). Embeddings that are closer to a tree-like structure (low $\delta$ values) capture the hierarchical relationships inherent in health-related data.
% The results from our experimental evaluation using the PDB186, Lung cancer, and Breast cancer datasets highlight the potential of incorporating $\delta$-hyperbolicity, ultrametricity, and Neighbor Joining (NJ) as a criterion for selecting and fine-tuning LLMs. Additionally, our analysis of synthetic data across different geometric spaces provides further insights into the behavior of hyperbolicity in various contexts, reinforcing the importance of understanding the geometric properties of embeddings.
% Future work could extend this research by exploring other geometric measures, applying this approach to more diverse datasets, and integrating these findings into real-world health applications to validate the practical benefits of hyperbolic embeddings.

% \textcolor{blue}{SARWAN: Need to update conclusion based on new metrics and results.}





% \begin{table}[h!]
%       \centering
%       \resizebox{0.5\textwidth}{!}{
% \begin{tabular}{@{\extracolsep{6pt}}p{0.5cm}lp{0.5cm}p{0.5cm}p{0.6cm}p{0.7cm}p{0.7cm}p{0.5cm}p{1.5cm}}
%     \toprule
%         \multirow{2}{*}{Embed.} & \multirow{2}{*}{Algo.} & \multirow{2}{*}{Acc. $\uparrow$} & \multirow{2}{*}{Prec. $\uparrow$} & \multirow{2}{*}{Recall $\uparrow$} & \multirow{2}{1.2cm}{F1 (Weig.) $\uparrow$} & \multirow{2}{1.2cm}{F1 (Macro) $\uparrow$} & \multirow{2}{1cm}{ROC AUC $\uparrow$} & Train Time (sec.) $\downarrow$\\
%         \midrule \midrule
%         % \cmidrule{2-9}  

%       \multirow{7}{1.2cm}{SeqVec} 
% & SVM & 0.05145 & 0.04687 & 0.05145 & 0.05661 & 0.05359 & 0.04533 & 607.12516		\\
% & NB & 0.04072 & 0.19305 & 0.04072 & 0.04176 & 0.02918 & 0.03741 & 19.22395       \\
% & MLP & 0.04657 & 0.13719 & 0.04657 & 0.03850 & 0.05438 & 0.05214 & 274.25843     \\
% & KNN & 0.07210 & 0.17445 & 0.07210 & 0.11939 & 0.10607 & 0.04312 & 942.97360     \\
% & RF & 0.04263 & 0.03893 & 0.04263 & 0.04269 & 0.04266 & 0.04038 & 11.91560       \\
% & LR & 0.04657 & 0.04523 & 0.04657 & 0.04569 & 0.04905 & 0.04912 & 83.08275       \\
% & DT & 0.05357 & 0.05589 & 0.05357 & 0.05494 & 0.05666 & 0.05631 & 12.44209       \\

% \midrule
% \multirow{7}{1.2cm}{ESM2}
% & SVM & 0.07886 & 0.08280 & 0.07886 & 0.09311 & 0.08184 & 0.06427 & 0.00469	\\
% & NB & 0.07076 & 0.06351 & 0.07076 & 0.07992 & 0.06390 & 0.04803 & 0.00062    \\
% & MLP & 0.07427 & 0.08397 & 0.07427 & 0.09030 & 0.07976 & 0.06248 & 0.14390   \\
% & KNN & 0.05297 & 0.05176 & 0.05297 & 0.05097 & 0.04247 & 0.03733 & 0.02459   \\
% & RF & 0.06389 & 0.06874 & 0.06389 & 0.08255 & 0.07217 & 0.05582 & 0.02241    \\
% & LR & 0.02187 & 0.09582 & 0.02187 & 0.02848 & 0.04005 & 0.04807 & 0.00466    \\
% & DT & 0.08527 & 0.09933 & 0.08527 & 0.10153 & 0.08945 & 0.07334 & 0.00402    \\

 
% \midrule
% \multirow{7}{1.2cm}{TAPE}
% & SVM & 0.09313 & 0.08535 & 0.09313 & 0.09084 & 0.09520 & 0.09321 & 0.01798	\\
% & NB & 0.08803 & 0.08360 & 0.08803 & 0.08822 & 0.09011 & 0.08765 & 0.00094    \\
% & MLP & 0.12690 & 0.10441 & 0.12690 & 0.12648 & 0.12681 & 0.11557 & 11.02384  \\
% & KNN & 0.03244 & 0.04184 & 0.03244 & 0.03122 & 0.03618 & 0.03912 & 0.38569   \\
% & RF & 0.07639 & 0.06249 & 0.07639 & 0.07393 & 0.07547 & 0.06681 & 0.04754    \\
% & LR & 0.09567 & 0.09223 & 0.09567 & 0.09449 & 0.09702 & 0.09590 & 0.05664    \\
% & DT & 0.07254 & 0.06236 & 0.07254 & 0.07031 & 0.07119 & 0.06745 & 0.01948    \\

% \midrule

% \multirow{7}{1.2cm}{ProtT5}
% & SVM & 0.04553 & 0.05350 & 0.04553 & 0.04448 & 0.04572 & 0.04808 & 0.01950	\\
% & NB & 0.06003 & 0.07512 & 0.06003 & 0.06450 & 0.07163 & 0.07797 & 0.00059    \\
% & MLP & 0.05841 & 0.06107 & 0.05841 & 0.05967 & 0.06455 & 0.06474 & 10.94558  \\
% & KNN & 0.04792 & 0.04382 & 0.04792 & 0.04942 & 0.04662 & 0.04474 & 0.30793   \\
% & RF & 0.03703 & 0.04341 & 0.03703 & 0.03714 & 0.03957 & 0.03956 & 0.03826    \\
% & LR & 0.05562 & 0.05406 & 0.05562 & 0.05485 & 0.05649 & 0.05590 & 0.01001    \\
% & DT & 0.07639 & 0.09442 & 0.07639 & 0.07657 & 0.07633 & 0.08460 & 0.02324    \\
     
%     \bottomrule
%   \end{tabular}  
%     }
%     \caption{Standard deviation results for different models and algorithms for \textbf{PDB186} dataset.}
%     \label{tbl_all_stf_data_results_pdb186}
% \end{table}

% \begin{table}[h!]
%       \centering
%       \resizebox{0.5\textwidth}{!}{
% \begin{tabular}{@{\extracolsep{6pt}}p{0.5cm}lp{0.5cm}p{0.5cm}p{0.6cm}p{0.7cm}p{0.7cm}p{0.5cm}p{1.5cm}}
%     \toprule
%         \multirow{2}{*}{Embed.} & \multirow{2}{*}{Algo.} & \multirow{2}{*}{Acc. $\uparrow$} & \multirow{2}{*}{Prec. $\uparrow$} & \multirow{2}{*}{Recall $\uparrow$} & \multirow{2}{1.2cm}{F1 (Weig.) $\uparrow$} & \multirow{2}{1.2cm}{F1 (Macro) $\uparrow$} & \multirow{2}{1cm}{ROC AUC $\uparrow$} & Train Time (sec.) $\downarrow$\\
%         \midrule \midrule

     
% \multirow{7}{1.2cm}{SeqVec}
% & SVM & 0.00850 & 0.01202 & 0.00850 & 0.00980 & 0.06036 & 0.03375 & 6.28171		\\
% & NB & 0.27801 & 0.04480 & 0.27801 & 0.29335 & 0.07373 & 0.04863 & 0.30832      \\
% & MLP & 0.01796 & 0.02203 & 0.01796 & 0.01836 & 0.07088 & 0.04319 & 98.55081    \\
% & KNN & 0.00756 & 0.01112 & 0.00756 & 0.01107 & 0.02505 & 0.02223 & 1.03840     \\
% & RF & 0.00481 & 0.00827 & 0.00481 & 0.00621 & 0.05260 & 0.03321 & 0.75593      \\
% & LR & 0.05222 & 0.00885 & 0.05222 & 0.03898 & 0.04887 & 0.03256 & 254.57333    \\
% & DT & 0.00710 & 0.00709 & 0.00710 & 0.00461 & 0.03187 & 0.03059 & 5.32662      \\

% \midrule
% \multirow{7}{1.2cm}{ESM2}
% & SVM & 0.01320 & 0.01663 & 0.01320 & 0.01424 & 0.04768 & 0.02114 & 4.68913		\\
% & NB & 0.01856 & 0.01490 & 0.01856 & 0.01507 & 0.02611 & 0.02914 & 0.09660      \\
% & MLP & 0.01739 & 0.01496 & 0.01739 & 0.01686 & 0.06158 & 0.03997 & 2.37537     \\
% & KNN & 0.00990 & 0.01774 & 0.00990 & 0.01283 & 0.06246 & 0.03353 & 0.24810     \\
% & RF & 0.01024 & 0.01633 & 0.01024 & 0.01281 & 0.05465 & 0.01972 & 0.21527      \\
% & LR & 0.00941 & 0.00857 & 0.00941 & 0.00721 & 0.01487 & 0.01736 & 27.16001     \\
% & DT & 0.02126 & 0.02629 & 0.02126 & 0.02224 & 0.05766 & 0.03063 & 1.38173      \\


% \midrule
% \multirow{7}{1.2cm}{TAPE}
% & SVM & 0.01346 & 0.01584 & 0.01346 & 0.01411 & 0.06311 & 0.03745 & 0.00127		\\
% & NB & 0.00997 & 0.01197 & 0.00997 & 0.00990 & 0.04767 & 0.03132 & 0.00138      \\
% & MLP & 0.01235 & 0.01432 & 0.01235 & 0.01350 & 0.02559 & 0.02373 & 0.35549     \\
% & KNN & 0.01024 & 0.01635 & 0.01024 & 0.01185 & 0.05105 & 0.02608 & 0.01132     \\
% & RF & 0.01149 & 0.01603 & 0.01149 & 0.01358 & 0.02507 & 0.02155 & 0.07465      \\
% & LR & 0.01420 & 0.01603 & 0.01420 & 0.01519 & 0.06293 & 0.04261 & 0.03193      \\
% & DT & 0.01723 & 0.01176 & 0.01723 & 0.01438 & 0.03369 & 0.02150 & 0.09068      \\

% \midrule
% \multirow{7}{1.2cm}{ProtT5}
% & SVM & 0.01346 & 0.01260 & 0.01346 & 0.01380 & 0.05012 & 0.01748 & 0.00532		\\
% & NB & 0.02434 & 0.02313 & 0.02434 & 0.02388 & 0.08211 & 0.04440 & 0.00281      \\
% & MLP & 0.01925 & 0.01172 & 0.01925 & 0.01613 & 0.04616 & 0.02035 & 1.54843     \\
% & KNN & 0.01224 & 0.01374 & 0.01224 & 0.01142 & 0.05051 & 0.02582 & 0.07756     \\
% & RF & 0.01781 & 0.02434 & 0.01781 & 0.01738 & 0.04168 & 0.02453 & 0.02489      \\
% & LR & 0.02239 & 0.02109 & 0.02239 & 0.02173 & 0.08363 & 0.03552 & 0.13927      \\
% & DT & 0.01371 & 0.01754 & 0.01371 & 0.01520 & 0.01058 & 0.01303 & 0.06655      \\
     
%     \bottomrule
%   \end{tabular}  
%     }
%     \caption{Standard deviation results for different models and algorithms for \textbf{Lung Cancer} dataset.}
%     \label{tbl_all_std_data_results_LC}
% \end{table}



% \begin{table}[h!]
%       \centering
%       \resizebox{0.5\textwidth}{!}{
% \begin{tabular}{@{\extracolsep{6pt}}p{0.5cm}lp{0.5cm}p{0.5cm}p{0.6cm}p{0.7cm}p{0.7cm}p{0.5cm}p{1.5cm}}
%     \toprule
%         \multirow{2}{*}{Embed.} & \multirow{2}{*}{Algo.} & \multirow{2}{*}{Acc. $\uparrow$} & \multirow{2}{*}{Prec. $\uparrow$} & \multirow{2}{*}{Recall $\uparrow$} & \multirow{2}{1.2cm}{F1 (Weig.) $\uparrow$} & \multirow{2}{1.2cm}{F1 (Macro) $\uparrow$} & \multirow{2}{1cm}{ROC AUC $\uparrow$} & Train Time (sec.) $\downarrow$\\
%         \midrule \midrule
%         % \cmidrule{2-9}  
    
% \multirow{7}{1.2cm}{SeqVec}
% & SVM & 0.00948 & 0.01335 & 0.00948 & 0.01137 & 0.02177 & 0.00973 & 3802.87105	\\
% & NB & 0.23678 & 0.04413 & 0.23678 & 0.26288 & 0.08182 & 0.04810 & 1.16807      \\
% & MLP & 0.01304 & 0.01768 & 0.01304 & 0.01360 & 0.04238 & 0.02531 & 335.81955   \\
% & KNN & 0.09328 & 0.02945 & 0.09328 & 0.06947 & 0.05295 & 0.02757 & 228.76138   \\
% & RF & 0.01174 & 0.01807 & 0.01174 & 0.01313 & 0.02157 & 0.01180 & 4.19326      \\
% & LR & 0.04008 & 0.01659 & 0.04008 & 0.02538 & 0.03783 & 0.04226 & 1254.46252   \\
% & DT & 0.01616 & 0.01889 & 0.01616 & 0.01551 & 0.01374 & 0.00414 & 12.45272     \\

% \midrule
% \multirow{7}{1.2cm}{ESM2}
% & SVM & 0.01447 & 0.02069 & 0.01447 & 0.01668 & 0.03476 & 0.05892 & 1.31803		\\
% & NB & 0.02562 & 0.00888 & 0.02562 & 0.02185 & 0.04030 & 0.02812 & 0.01634      \\
% & MLP & 0.02099 & 0.02567 & 0.02099 & 0.02289 & 0.05487 & 0.02629 & 5.96107     \\
% & KNN & 0.02169 & 0.02387 & 0.02169 & 0.02322 & 0.06704 & 0.06635 & 0.18274     \\
% & RF & 0.01976 & 0.02739 & 0.01976 & 0.02322 & 0.03967 & 0.02853 & 0.33201      \\
% & LR & 0.01005 & 0.02292 & 0.01005 & 0.01449 & 0.02354 & 0.05516 & 8.20804      \\
% & DT & 0.02191 & 0.02509 & 0.02191 & 0.02357 & 0.07411 & 0.06816 & 0.38781      \\

% \midrule
% \multirow{7}{1.2cm}{TAPE}
% & SVM & 0.01313 & 0.01704 & 0.01313 & 0.01464 & 0.05405 & 0.03009 & 0.02476		\\
% & NB & 0.01765 & 0.02128 & 0.01765 & 0.01967 & 0.03097 & 0.03209 & 0.00430      \\
% & MLP & 0.01944 & 0.01815 & 0.01944 & 0.01890 & 0.05160 & 0.02974 & 51.26331    \\
% & KNN & 0.01180 & 0.01491 & 0.01180 & 0.01331 & 0.03447 & 0.02721 & 0.42028     \\
% & RF & 0.02320 & 0.03875 & 0.02320 & 0.03043 & 0.10826 & 0.05426 & 0.27117      \\
% & LR & 0.00577 & 0.00437 & 0.00577 & 0.00532 & 0.02881 & 0.02060 & 0.24390      \\
% & DT & 0.01890 & 0.01284 & 0.01890 & 0.01480 & 0.06272 & 0.03327 & 0.23943      \\

% \midrule
% \multirow{7}{1.2cm}{ProtT5}
% & SVM & 0.01200 & 0.01721 & 0.01200 & 0.01413 & 0.04908 & 0.04241 & 0.00656		\\
% & NB & 0.02367 & 0.02542 & 0.02367 & 0.02412 & 0.05954 & 0.03831 & 0.00609      \\
% & MLP & 0.02025 & 0.02171 & 0.02025 & 0.02073 & 0.06092 & 0.05220 & 1.21485     \\
% & KNN & 0.01816 & 0.02376 & 0.01816 & 0.02114 & 0.02956 & 0.01631 & 0.01540     \\
% & RF & 0.01890 & 0.02547 & 0.01890 & 0.02232 & 0.08446 & 0.03270 & 0.03397      \\
% & LR & 0.01289 & 0.01798 & 0.01289 & 0.01514 & 0.04249 & 0.04834 & 0.04502      \\
% & DT & 0.03063 & 0.03345 & 0.03063 & 0.03145 & 0.05199 & 0.04237 & 0.08053      \\
%     \bottomrule
%   \end{tabular}  
%     }
%     \caption{Standard deviation results for different models and algorithms for \textbf{Breast Cancer} dataset.}
%     \label{tbl_all_std_data_results_BC}
% \end{table}

% \clearpage
% \appendix

% \section*{Data Statistics}



% https://docs.google.com/document/d/1CF-wSjN-cBg3t0BJT84hGlSJgVhfopv-iwYsPQPet8o/edit?usp=sharing

% \clearpage
% \appendix

% \section{Proposed Approach}
% \label{supp_proposed_approach}




% \end{document}


\bibliographystyle{plainnat-revised}
% \bibliographystyle{IEEEtran}
\bibliography{references}


\end{document}